\documentclass[11pt]{article}
\usepackage[T1]{fontenc}
\usepackage{lmodern}
\usepackage{needspace}
\usepackage[margin=0.85in]{geometry}
\PassOptionsToPackage{obeyspaces,spaces}{url}
\usepackage{amsmath,amssymb,mathtools,booktabs,longtable,array,enumitem,microtype,xcolor,xurl,hyperref,fancyhdr}
\hypersetup{colorlinks=true,linkcolor=blue!45!black,urlcolor=blue!45!black,
 pdftitle={Perturbative matching in the November 4 OpenAI Yang-Mills construction},
 pdfauthor={Working notes for the Yang-Mills group},
 pdfsubject={Perturbative matching: proof summary, definition glossary, supplied tests and independent test specifications}}
\setlength{\parindent}{0pt}
\setlength{\parskip}{5pt plus 1pt}
\setlength{\headheight}{14pt}
\setlist{leftmargin=1.5em,itemsep=2pt,topsep=3pt}
\pagestyle{fancy}\fancyhf{}
\fancyhead[L]{\small November Yang--Mills: perturbative matching}
\fancyhead[R]{\small Working draft, 13 November 2026}
\fancyfoot[C]{\thepage}
\newcommand{\R}{\mathbb R}
\newcommand{\C}{\mathbb C}
\newcommand{\N}{\mathbb N}
\newcommand{\MS}{\overline{\mathrm{MS}}}
\newcommand{\code}[1]{{\small\nolinkurl{#1}}}
\newcommand{\src}[1]{\textsf{[#1]}}
\newcommand{\Conf}{\operatorname{Conf}}
\newcommand{\cO}{\mathcal O}
\makeatletter
\newcommand{\compactcontents}{\begingroup\small\setlength{\parskip}{0pt}
\renewcommand{\l@section}{\@dottedtocline{1}{0em}{1.8em}}
\tableofcontents\endgroup}
\makeatother
\begin{document}
\begin{center}
{\LARGE\bfseries Perturbative matching in the November 4\\[3pt]
OpenAI Yang--Mills construction}\par
\vspace{7pt}
{\large What is compared, how the comparison is proved,\\
and what remains to validate}\par
Working notes for the Yang--Mills group\quad\textbullet\quad 13 November 2026
\end{center}

\textbf{Purpose.} This is a first account of the November submission's
perturbative-matching theorem, separate from our comparison of the October
and November existence results. It explains the actual objects in the
statement and traces the main proof connections in the Lean source.
The intended reader knows quantum field theory but need not know this
formalization's terminology. The appendix gives a glossary of the principal formal definitions.
Sections~\ref{sec:inventory} and~\ref{sec:independent} list the supplied
tests and our independent test specifications, with their current status.

\textbf{Main finding.} The submitted target requires much more than a
free-Maxwell ultraviolet limit. A single constructed Wilson theory has,
in prescribed scale-dependent field coordinates, asymptotic expansions
through every finite order for separated correlations and the full
finite-degree OPE output. Their coefficients are fixed by a continuum
graph definition and a specified change of renormalization scheme.
The estimates include position derivatives and apply to the mixed
hierarchy and every character phase where the correlation statement
quantifies over phases. An independently characterized stress-tensor
series is matched as well.

\textbf{Review status.} The Matching target has already passed the two
recorded Comparator runs described in Section~\ref{sec:receipts}. This
document records a source investigation, not a completed independent
validation of every graph rule, subtraction convention, or analytic
estimate. In particular, ``the reference is defined independently'' and
``the reference is the intended standard Yang--Mills perturbation theory''
are distinct claims. We have traced an all-order identification between
the explicit graph sums and a separately specified renormalization-flow
problem, as well as the connection of its seeds to the classical action.
Section~\ref{sec:graph-audit} explains this evidence and the remaining
checks of conventions. Additional Lean checks establish the free gauge
line, both ghost-line orientations, and the cubic ghost--gluon vertex
in explicit conventions. The full finite ghost-loop benchmark is now
proved in Lean and has passed a separate Comparator run. It identifies
the actual submitted, regulator-removed graph in the infinite infrared
cutoff limit with the explicit $\MS$ answer, including both finite
constants, for all nonzero momenta, tensor components, colors and positive
renormalization scales. Appendix~\ref{app:finite-ghost} explains the
calculation and its formal proof. This is an independent check of one
gauge-dependent graph contribution; it does not identify every
renormalized amplitude with a conventional calculation.

\setcounter{tocdepth}{1}
\compactcontents
\clearpage

\section{The final theorem and its quantifiers}
\label{sec:statement}

The public challenge asks for
\begin{quote}
\code{YangMills.perturbativeMatching (g : GaugeGroup) (simple : g.IsSimple)}\par
\code{:}\ $\exists$\ \code{run : Wilson.Construction g, Nonempty (MatchedRun run)}.
\end{quote}
The submitted solution supplies \code{WilsonProof.publishedRun} and
\code{WilsonProof.matchedRun}. The object \code{MatchedRun} contains one
perturbative reference, one microscopic Wilson family with its native
insertion coordinates, one family of finite-period continuum kernels,
and the matching data for its chosen large-volume refinements. The
reference is chosen before those volume refinements. \src{T}

For each strictly increasing requested sequence of volume indices, the
construction selects a subsequence and a limiting hierarchy. A
\code{MatchedRefinement} ties the matching estimates to that same
hierarchy, its lattice limit, its phase space, and its physical OPE.
Its field frame and coupling are common to all arities, perturbative
orders, input fields, output cutoffs, and phases within that refinement.
One does not choose a different theory or a fresh normalization to make
each coefficient agree.

There is also a proved implication
\begin{quote}
\code{MatchedRun.realizes : run.Realizes (maxwellReference g)}.
\end{quote}
Thus the matching witness satisfies November's existence requirements
for the same run. This is stronger than placing unrelated existence and
matching assertions next to one another. It does not assert equality
of every witness that could separately be chosen by the existence proof.

The large-volume quantifier remains subsequential. Neither
\code{MatchedRun} nor the uniqueness of its perturbative reference proves
that all large-volume limits are equal. The reference is common across
the refinements of the run; the measured physical frames and coupling
are supplied on each limiting hierarchy. \src{T,V}

\section{The fields, the radius, and the two sides of the comparison}
\label{sec:objects}

Let $B=\texttt{maxwellReference}(g)$ and write $\mathcal V_{\le d}$ for
\code{B.sector d}. This is the finite-dimensional space of gauge-invariant
free on-shell curvature-jet polynomials through engineering degree $d$.
It is a filtered space: membership in $\mathcal V_{\le d}$ need not mean
that a field is homogeneous of degree exactly $d$. A derivative of order
$j$ of curvature has weight $j+2$. The same polynomial has compatible
representations at all larger cutoffs. The physical hierarchy $\mathcal H$
has corresponding spaces \code{H.Field d}. \src{F,N}

The matching data provide linear isomorphisms
\begin{equation}\label{eq:frame}
 J_r^{(d)}:\mathcal V_{\le d}\longrightarrow\mathcal H.\mathrm{Field}(d),
 \qquad r>0,
\end{equation}
and a measured real coupling $g_X(r)$. Here $r$ is a \emph{physical
short-distance radius in the continuum theory}, not the lattice spacing
or a volume index. The frame includes the rescaling and field
renormalization needed for a finite ultraviolet limit. Consequently one
must not insert an additional factor $r^{\sum d_i}$ into the public
matching formula below. For homogeneous unrenormalized fields, such
engineering factors are part of what the frame implements.

For $v_i\in\mathcal V_{\le d_i}$, a center $o\in\R^4$, and a separated
shape $x=(x_1,\ldots,x_n)$, define
\begin{equation}\label{eq:measured}
 S_{\sigma,r}(v;x,o)=
 S_\sigma\bigl(J_r^{(d_1)}v_1,\ldots,J_r^{(d_n)}v_n;
                   o+rx_1,\ldots,o+rx_n\bigr).
\end{equation}
The index $\sigma$ is either the mixed hierarchy (Lean's \code{none}) or
a character phase (\code{some character}). Both the positions and the
inserted fields change with $r$. The operator spaces and their relation
to the fixed lattice family remain the same.

On the other side are numerical functions
$C^{\MS}_{n,k}(v;x)$, defined by graphs from the gauge group and the free
polynomial labels. After the prescribed field and coupling conversion,
they become $C^X_{n,k}(v;x)$. These reference functions have no phase
argument and no dependence on the selected Wilson hierarchy. Their
dependence on $g$ includes the Lie bracket and color contractions, not
only the dimension of the Lie algebra. \src{G,N}

Four limits must be kept distinct. The physical construction removes
the lattice ultraviolet cutoff at fixed volume, then takes a selected
large-volume limit. The matching theorem then sends $r\downarrow0$
within that continuum hierarchy. The graph definition has its own
dimensional regulator and infrared cutoff, whose removal is a separate
perturbative calculation (Section~\ref{sec:graphs}).

\section{How the perturbative reference is defined}
\label{sec:graphs}

\subsection{Local vertices, propagators, and graph assembly}

The public definition \code{graphCorrelation} first includes all input
fields in a common finite cutoff, expands them in its graded basis,
evaluates each basis tuple with \code{YM.MS.Graphical.fullGraph}, and
extends multilinearly. The local insertion polynomials come from
\code{OpenAxis.Plot}; the common basis is \code{Press.portDeck}. This
definition does not take a continuum Wilson correlation as an argument.
\src{G}

The graph machinery represents families of scale-ordered graphs by
three operations: a local vertex, joining two connected components by
an edge, and closing a pair of legs inside a component. A \code{Diagram}
is therefore not a single fully labelled incidence graph. Its evaluator
performs the field-index sums, graded permutations, contractions and
scale integrations. The enumeration includes repeated closures and
recursively allocates positive coupling order to bulk vertices.
\src{D}

The bulk input is explicitly
\[
 \chi=g B_1+g^2B_2.
\]
In \code{Motif/Receipt.lean}, $B_1$ contains the cubic gauge vertex
\code{top3} and the term \code{cape (blazeV)} in the ghost/source
formalism; $B_2$ is the quartic gauge vertex \code{top4}. The displayed
gauge terms carry $-1/4$ times, respectively, the two cross terms
$F_0[A,A]+[A,A]F_0$ and the term $[A,A][A,A]$. The sign must be read with
the source's exponent and anti-Hermitian trace conventions. These
formulas are evidence for the intended Yang--Mills vertices; a human
validation should still follow their tensor and graded-sign conventions
down to the kernels.

The propagator operations use \code{Landau.choir} and heat-shell kernels.
The edge rule includes graded signs for field parity. The distinction
between a bridge and a loop closure is explicit: only closure increments
the dimensionally continued loop count. This is a scale-integral
organization of the calculation, rather than a list of momentum-space
Feynman integrals in familiar notation. A reader should verify its
equivalence to the intended gauge-fixed rules, rather than infer it from
the names of the modules. \src{D}

\subsection{Subtraction and removal of regulators}

At a graph node, \code{integrateRen} first integrates the shell expression,
including the counterterms of its children. It then adds the negative
Laurent principal part of the resulting \emph{summed} sector. The
counterterm is defined by finite Taylor-coefficient extraction after
multiplication by a sufficient power of the dimensional parameter.
The code therefore specifies a subtraction operation; it does not leave
the renormalized answer as an arbitrary parameter. \src{D}

\code{regularPart} takes the dimensional parameter to zero through
nonzero complex values after the counterterm expressions have been
assembled. \code{irRemoved} subsequently takes the large infrared-length
limit. The remaining moment assembly combines connected cluster
contributions into full moments. The public correlation definition
uses the logarithmic reference scale
\begin{equation}\label{eq:ms-scale}
 y_{\MS}=\log 1+\tfrac12\bigl(\gamma_E-\log(4\pi)\bigr).
\end{equation}
Thus the $\MS$ convention is represented by an actual constant in the
definition. \src{G,D}

There are accompanying theorems, not only totalized limit definitions.
\code{WilsonGraphNamedGraphForestLimits.physical} states, for separated
positions at every fixed perturbative order, convergence of the
dimensionally regulated expression at each positive infrared length,
followed by convergence as that length tends to infinity to the issued
graph coefficient. \code{WilsonGraphPublicForestLimits.convertedLimits}
transports these limits through the finite scheme conversion.
This is a sequential removal of regulators; it should not be described
as an unrestricted joint limit. \src{R}

The present investigation has followed these definitions and theorems
into the matching argument. It has not independently checked every
graph multiplicity or every estimate used to remove the regulators.

\subsection{Why these graphs are more than a stipulated reference}
\label{sec:graph-audit}

There is a substantial proof between the graph definitions and the
matching theorem. The submission does not merely name an array of
numbers ``Yang--Mills perturbation theory'' and compare against it.
It specifies a flow equation, its classical input, analytic conditions,
and pole-subtraction boundary conditions; proves uniqueness; and proves
that the explicit graph sums solve that problem. This shifts the main
validation question to a concrete one: do these input data and conditions
express the intended perturbation theory? \src{B}

Flow equations are an established way to formulate perturbative
Yang--Mills theory. For relevant mathematical comparisons, see
Efremov, Guida and Kopper on perturbative renormalization of SU(2),
and Fr\"ob and Holland on all-order perturbative OPEs.\footnote{
\href{https://arxiv.org/abs/1704.06799}{A.~N.~Efremov, R.~Guida and C.~Kopper,
\emph{Renormalization of SU(2) Yang--Mills theory with Flow Equations}}
(2017);
\href{https://arxiv.org/abs/1603.08012}{M.~B.~Fr\"ob and J.~Holland,
\emph{All-order existence of and recursion relations for the operator
product expansion in Yang--Mills theory}} (2016).}
These are references for the method and possible comparison targets.
We have not established that the submitted code follows either paper,
or that its regulator and renormalization conventions coincide with theirs.

First consider the action. In the graded polynomial algebra the
curvature and BRST variations are explicitly
\[
 F_{\mu\nu}(g)=\partial_\mu A_\nu-\partial_\nu A_\mu
                 +g[A_\mu,A_\nu],\qquad
 sA_\mu=\partial_\mu c+g[A_\mu,c],\qquad
 sc=-\frac g2[c,c].
\]
The object \code{Motif.ploy} combines the curvature-squared log density
with source terms for these variations. Its gauge part is
$-\frac14\sum_{\mu,\nu,a}(F^a_{\mu\nu})^2$ in the source's component
conventions. In particular this is a log-density convention, so the
minus sign should not be silently changed to the sign of a Euclidean
action in an $e^{-S}$ convention. The theorem \code{Motif.Class_at} proves
\begin{equation}\label{eq:classical-input}
 gB_1+g^2 B_2=\mathrm{ploy}(g)-\mathrm{ploy}(0).
\end{equation}
Thus the seeds are proved to be the interaction part of that specified
classical expression. \code{Motif.AllMaster} proves its classical master
identity on momentum-conserving words. This is a classical identity;
it should not be presented as, by itself, a proof of every quantum Ward
identity. \code{Motif.kernel} separately computes the free gauge Hessian
as $-\delta^{ab}(\delta_{\mu\nu}p^2-p_\mu p_\nu)$. \src{B}

The gauge line has transverse Landau form. Our small additional check
in \code{DiagramReferenceAudit.lean} proves, directly from the submitted
line definition and for $p^2\ne0$, that
\begin{equation}\label{eq:gauge-line-audit}
 2\,\mathrm{halfCov}^{ab}_{\mu\nu}(s,p)=
 \delta^{ab}\,[1-H(e^{2s}p^2)]
 \left(\frac{\delta_{\mu\nu}}{p^2}
                   -\frac{p_\mu p_\nu}{(p^2)^2}\right),
 \qquad H(u)=(1+u)e^{-u}.
\end{equation}
The factor two matters: the code uses half-lines for ordered
contractions. The submitted theorem \code{length_recovers_line}
then identifies a scale integral of the shell line with the difference
of these filtered propagators. Our check concerns the free gauge line
and its color delta. It does not yet prove the position-space curvature
comparison in test I1 or validate an interacting coefficient.

\paragraph{The ghost and source dictionary.}
The line table uses a scalar odd ghost $c$ and an odd vector source
$K_\mu$, called \code{anti} in the code. Their ghost numbers are $+1$
and $-1$ and their engineering dimensions are one and two. A further
even source $c^*$ has ghost number $-2$ and no propagator. These are
BRST source variables; they should not all be interpreted as physical
fields. The source expression is proved in \code{Halo.Surf.reach} to be
\begin{equation}\label{eq:ghost-source}
 \mathrm{cape}(s)=\sum_{a,\mu}K^a_\mu(D_\mu c)^a
                         -\sum_a c^{*a}(sc)^a.
\end{equation}
This displays the connection to a scalar antighost. Choose Fourier
derivatives $\partial_\mu\mapsto ip_\mu$, ordered free contractions
$\langle c^a(p)\bar c^b(-p)\rangle=\delta^{ab}/p^2$ and
$\langle\bar c^a(p)c^b(-p)\rangle=-\delta^{ab}/p^2$, and substitute
\begin{equation}\label{eq:antighost-dictionary}
 K^a_\mu=-\partial_\mu\bar c^a,\qquad c^{*a}=0.
\end{equation}
The two ghost orientations then give
\[
 \langle c(p)K_\mu(-p)\rangle
   =\frac{1}{p^2}(+ip_\mu),\qquad
 \langle K_\mu(p)c(-p)\rangle
   =(-ip_\mu)\left(-\frac{1}{p^2}\right).
\]
Both equal $ip_\mu/p^2$, as in the submitted line table; color deltas
are suppressed here. The sign in the first formula comes from the
momentum on the second leg; the sign in the second is the exchange of
two odd fields. These are identities of symbols away from zero momentum,
not prescriptions for extending singular distributions.

Under the same substitution, equation~\eqref{eq:ghost-source} becomes
$-\sum_{a,\mu}(\partial_\mu\bar c^a)(D_\mu c)^a$. This is the negative
of the usual ghost action in an $e^{-S}$ convention, consistent with
the log-density convention of \code{ploy}. Our
\code{GhostConventionCheck.lean} proves both uncolored propagator
identities directly from the submitted definitions. It also proves that
the full classical log density has the derivative-normalized vertex
\[
 (K^a_\mu,A^b_\nu,c^c):\quad
       g\delta_{\mu\nu}f^{bc}{}_{a},
\]
and that substituting~\eqref{eq:antighost-dictionary} gives
\begin{equation}\label{eq:ghost-vertex-audit}
 (\bar c^a(q),A^b_\nu(p),c^c(k)):\quad
       -igq_\nu f^{bc}{}_{a}.
\end{equation}
Here $f^{bc}{}_{a}$ means precisely the submitted bracket component
\code{rune b c a}. No additional factorial is inserted: the rows are
derivatives in the displayed external-field order. Momentum conservation
is imposed separately when forming an amplitude. \src{K}

This resolves the free-line and local cubic-vertex signs in an explicit
scalar-antighost convention. It is not an all-order theorem that the
renormalized graph systems agree under gauge fixing. In particular,
the substitution makes $K$ longitudinal and is not an invertible map
from the unrestricted source space. Source insertions, the BRST
identities used for composite fields, and their counterterms still
require a compatible dictionary.

\subsection{The all-order recursion and its uniqueness}

The bulk flow uses two operations: a linear contraction that closes a
pair of legs, and a bilinear contraction that joins two components.
Schematically its coefficient equation has the form
\[
 \mathcal D X_n=\mathcal L X_n+
       \frac12\sum_{\substack{a+b=n\\a,b>0}}\mathcal B(X_a,X_b).
\]
Here $\mathcal D$ includes the scale and homogeneity bookkeeping,
and $\mathcal L,\mathcal B$ include the shell kernels and graded signs;
this display is a description of the recursion, not a replacement for
its full definition. For one marked insertion the bilinear term is
$\sum_{a+b=n,\,a>0}\mathcal B(X_a,Y_b)$, without the bulk
component-order factor $1/2$. Joining components therefore reduces
coupling order in the unknown factor. Repeated closure is controlled
by the finite leg-degree bound. These two reductions organize the
induction. \src{B}

The Lean predicate \code{Dock.Joint} includes the flow equation,
meromorphic dependence on the regulator, regularity after the
appropriate assembly, pure-pole counterterms, their homogeneity,
and a prescribed ultraviolet boundary value. \code{Bob} adds the
closure term and degree bounds. \code{leader} and \code{comet}
impose the resulting coefficient equations on the bulk and marked
series. \code{MS_Data} is the data of such solutions; existence is
proved in \code{bookedMS}. Uniqueness is proved by
\code{Dock.uniqueK}, \code{unique_B}, and \code{unique_M}.

On the graph side, \code{bridgeForce_eq} identifies graph joining with
the bilinear flow operation. \code{closeForce_eq} identifies graph
closure with the linear operation in the sense of regulator germs.
\code{ambient_bob} and \code{mark_bob} sum the actual enumerated graph
families and prove the required equations, including their weights.
The final consequences are
\code{Diagram.income_bulk} and \code{Diagram.income_mark}:
the graph counterterms equal those of any \code{MS_Data} with the same
input, and every coefficient page agrees as a germ near the punctured
regulator origin. This last qualification matters: it is not an
assertion that arbitrary representatives agree at the singular point
$z=0$. \src{B}

This argument is evidence about multiplicities at every order, not just
a check of a few low-order diagrams. A weighted graph sum satisfying
the uniquely specified recursion can establish the correct aggregate
coefficient without a separate bijection to an incidence-graph list.
Its physical force still depends on validating the contraction
operations and boundary conditions themselves.

\subsection{Subtraction, dimensional continuation, and the final bridge}

The graph subtraction also has proved properties. \code{block}
shows that \code{integrateRen} solves the required response problem;
\code{ct_polynomial} shows that its counterterm equals its local Taylor
projection of the specified degree. The pole prescription is applied
to the assembled scale sector, with counterterm insertions already
included in its children. These facts justify calling it a local
pole-subtraction construction. They should not be silently upgraded
to a separately verified equality with every conventional
graph-by-graph forest prescription. \src{B}

There is also a useful normalization check on dimensional continuation.
\code{radialClose_integer} proves that the continued contraction at
$z=n$ agrees with ordinary integration in $4+2n$ dimensions, on its
stated admissible kernels and lifted external data.
\code{pitch_normalization} proves the measure factor
\[
 e^{-2yn}(2\pi)^{-(4+2n)}.
\]
These are genuine integral identities, including the Fourier and mass
normalizations. Agreement at the nonnegative integers alone would not
uniquely determine an analytic continuation near zero: those points
have no finite accumulation point. Consequently the explicit
noninteger-dimensional rule and its treatment of evanescent tensors
must also be reviewed. Terms vanishing at $z=0$ can affect finite
answers when multiplied by poles.

\paragraph{An actual pole insertion retains the finite regulator term.}
The source does address this last danger in a concrete example, beyond
the elementary algebra of multiplying a pole by a Taylor series.
In \code{ScaleForest.Bubbles}, \code{bubble G true true} is built from
two cubic interaction vertices and two ghost contractions.
\code{TwoPoint.poleK} identifies its counterterm with a simple pole;
\code{TwoPoint.ghost_bill} evaluates its local two-gauge-field tensor
with the submitted color trace and loop normalization. This is the
ghost-loop contribution to gauge-field self-energy. It is not a
two-point function of external ghosts. \src{S}

The subsequent calculation inserts this pole into the specified family
of graphs with two marked insertions. For admitted separated placements,
at retained infrared length, \code{Triangle.ghost_input} proves an
identity of regulator germs of the form
\begin{equation}\label{eq:actual-pole-complement}
 H(z)=B(z)+\frac{D(z)}{z},\qquad z=-\varepsilon.
\end{equation}
Here $B$ is the actual prepared graph contribution, $D$ is the
still-regulated surrounding graph with the calculated local pole tensor
inserted, and $H$ is the resulting renormalized contribution.
The pole sign is included in $D$; one must not change $z$ to
$\varepsilon$ without also changing that sign.
For any of the projected components in the theorem, write $d(z)$
for the corresponding component of $D(z)$. Then
\code{Triangle.premature} proves
\begin{equation}\label{eq:premature-pole}
 \lim_{z\to0}\left(B(z)+\frac{d(0)}{z}\right)
       =H(0)-d'(0).
\end{equation}
In this display $B$ and $H$ likewise denote the selected scalar
components. Thus freezing the complement before pole multiplication
loses the finite derivative term. This statement concerns the defined
graph integrals, not just a freely specified pair of Taylor coefficients.
\code{Triangle.DScaleJet} also computes the change of this derivative
under rescaling the placements and length. Neither result asserts that
$d'(0)$ is nonzero in every component or that its contribution cannot
cancel against other graphs in a physical amplitude.

The generic one-loop calculation in \code{ScaleForest.PageOne} makes
the same issue visible. If $T(z)$ is its divergent Taylor tensor, the
finite value contains
\[
 T'(0)-2(y+s)T(0)
\]
in addition to the regular integral, where $y$ and $s$ are its logarithmic
mass and length parameters. \code{simple_finite} proves that value and
\code{simple_pole_rule} identifies the pole with the prescribed unique
counterterm under the stated flow hypotheses. These results support the
claim that finite regulator terms are retained within the submitted
prescription. An independent comparison of that prescription with the
intended $\MS$ amplitudes remains necessary. \src{S}

\paragraph{An independent one-loop benchmark.}
The ghost-loop contribution to the gluon two-point vertex provides a
manageable next test. It is a gauge-dependent intermediate contribution,
not a complete gauge-invariant correlation function. In particular it
need not be transverse by itself. We strip off $g^2$, retain the color
factor $C^{ab}=C_A\delta^{ab}$, and define our sign by the Euclidean
dimensionally continued expression
\begin{equation}\label{eq:ghost-loop-integral}
 \mathcal B_{\mu\nu}^{ab}(p)
 =-C^{ab}\mu_{\rm DR}^{2\varepsilon}
   \int\frac{d^{4-2\varepsilon}k}{(2\pi)^{4-2\varepsilon}}
       \frac{k_\mu(k+p)_\nu}{k^2(k+p)^2},\qquad p^2>0.
\end{equation}
This fixes the sign before comparison with a counterterm. The minus
sign is the one obtained for the quadratic log-density contribution
from expanding the ghost determinant in the conventions above.
Use $\mu_{\rm DR}^2=\bar\mu^2 e^{\gamma_E}/(4\pi)$ and write
$\ell=\log(\bar\mu^2/p^2)$.

Here is the independent calculation to be checked. A Feynman parameter
$x$ and the shift $k\mapsto k-xp$ reduce the even numerator to
$k_\mu k_\nu-x(1-x)p_\mu p_\nu$. Rotational integration, followed by
the parameter integral, gives
\begin{align}
 \mathcal B_{\mu\nu}^{ab}(p;\varepsilon)
 &=\frac{C^{ab}}{16\pi^2}\frac{Q(\varepsilon,\ell)}{\varepsilon}
   \left(\frac{\delta_{\mu\nu}p^2}{2(1-\varepsilon)}
                         +p_\mu p_\nu\right),
 \label{eq:ghost-loop-gamma}\\
 Q(\varepsilon,\ell)
 &=e^{(\ell+\gamma_E)\varepsilon}\Gamma(1+\varepsilon)
       \frac{\Gamma(2-\varepsilon)^2}{\Gamma(4-2\varepsilon)}.
 \nonumber
\end{align}
The parameter integral at zero regulator is
$\int_0^1x(1-x)\,dx=1/6$. The first derivative of $Q$ is
$(\ell+5/3)/6$. Expanding the additional factor $1/(1-\varepsilon)$
then gives
\begin{equation}\label{eq:ghost-loop-benchmark}
 \mathcal B_{\mu\nu}^{ab}(p;\varepsilon)
 =\frac{C^{ab}}{16\pi^2}
 \left\{\frac{\delta_{\mu\nu}p^2}{12}
       \left(\frac1\varepsilon+\ell+\frac83\right)
    +\frac{p_\mu p_\nu}{6}
       \left(\frac1\varepsilon+\ell+\frac53\right)\right\}
       +O(\varepsilon).
\end{equation}
These two finite constants agree with Rodigast's explicit ghost-loop
calculation, equation~(4.30), after allowing for his Minkowski
conventions and regulator $4-d=2\varepsilon$.\footnote{
\href{https://qft.physik.hu-berlin.de/wp-content/uploads/2012/10/DA-andreas.pdf}{
A.~Rodigast, \emph{One-Loop Divergences of the Yang--Mills Theory Coupled
to Gravitation}}, diploma thesis, revised 20 August 2008,
Section~4.3.1, equation~(4.30), printed p.~27. The quoted expression is
the pure Yang--Mills ghost loop, without gravitational corrections.}

Our \code{GhostLoopBenchmark.lean} defines the Gamma expression without
using the submitted graph evaluator. It proves its recurrence reduction,
its value and derivative at zero, and the punctured-regulator limit
after subtracting the pole. Thus the constants $8/3$ and $5/3$ are
derived from this independent analytic expression in Lean, rather than
merely inserted as a proposed answer. The reduction of the dimensional
momentum integral to that Gamma expression is the human calculation
above; it has not been formalized in this check.

The same file proves \code{submitted_pole}: the actual submitted
\code{TwoPoint.poleTensor} agrees with the residue in
equation~\eqref{eq:ghost-loop-benchmark}, including
$C_A\delta^{ab}$. The source counterterm is $R/z$ with $z=-\varepsilon$,
so it contributes $-R/\varepsilon$ and cancels the displayed positive
bare pole. This checks the full two-index residue, not just a transverse
projection or a coefficient used to set the coupling. The color
comparison uses the submission's proved \code{sums_angles} identity;
it does not depend on a separately chosen SU(2) basis. \src{S}

\paragraph{The completed finite comparison.}
There is a proved connection to the finite graph, rather than only to its
pole. Our \code{submitted_finite_reading} applies the submission's
\code{twoVertex_laurent} theorem and identifies the actual
\code{Outlook.four} reading with \code{oneLoopFinite}. Write $W(z)$
for \code{TwoPoint.ratePair G true true}, $T(z)$ for its
\code{highTensor}, and $\mathcal R_s$ for the regular part, namely
\code{plow (-4) s (ordinary 4 W 0)}, evaluated on the chosen external
word. The exact formula is
\begin{equation}\label{eq:retained-ghost-finite}
 B_s(y;p)=\mathcal R_s(p)+T'(0;p)-2(y+s)T(0;p).
\end{equation}
Here $s=\log\rho$ is the retained infrared length and
$y=\log\bar\mu+(\gamma_E-\log(4\pi))/2$. The pole comparison already
identifies $T(0;p)=-C^{ab}R_{\mu\nu}(p)$, where
\[
 R_{\mu\nu}(p)=\frac1{16\pi^2}
        \left(\frac{\delta_{\mu\nu}p^2}{12}
                          +\frac{p_\mu p_\nu}{6}\right).
\]
Consequently \code{submitted_scale_coefficient} proves, for every
retained $s$,
\begin{equation}\label{eq:ghost-scale-difference}
 B_s(y';p)-B_s(y;p)=2(y'-y)C^{ab}R_{\mu\nu}(p).
\end{equation}
This is exactly the change of our independent finite reference when
$\bar\mu$ changes. It checks the coefficient of the renormalization-scale
logarithm. The following full comparison also checks the momentum
dependence and the finite constants.

The theorem \code{GhostSolution.acceptance} in
\code{checks/matching/GhostLoopSolution.lean} proves the exact predicate
\code{SubmittedFinitePartTest}:
apply the submission's regulator-removal reader \code{Outlook.four}
to this bubble, use its $\MS$ logarithmic scale, and let the infrared
length tend to infinity at fixed nonzero external momentum. The result
equals the finite part in~\eqref{eq:ghost-loop-benchmark}.
Our \code{finite_test_iff_reference} proves that this entire requirement
is equivalent to \code{ReferenceScaleTest}, which fixes $\bar\mu=1$
in the chosen units and still tests every nonzero momentum and both
tensor channels. The proof goes directly to the full target: it evaluates
the actual continued source force and its regulator derivative, solves
the scale equation with the submitted ultraviolet boundary condition,
and removes the infrared cutoff. Covariance and momentum rescaling
extend the unit-momentum calculation to every component and momentum.
No finite-part identity is assumed. A fresh build and axiom audit of
15 review modules passed, followed by Comparator on the frozen full
target on 13 November. The proof and the limits of that check are
described in Appendix~\ref{app:finite-ghost}.

\paragraph{An alternative route through the integrated propagators.}
The cutoff is explicit. \code{halfP_components} in
\code{ScaleForest/HeatScale.lean} identifies
the integrated line, including the Wick factor $1/2$, as the usual
free line multiplied by
\[
 f_\rho(q)=1-(1+\rho^2q^2)e^{-\rho^2q^2}.
\]
Our additional Lean lemmas prove $0\le f_\rho\le1$ and identify the
pointwise limit of that actual submitted line as $\rho\to\infty$
at nonzero $q$. This does not permit taking the limit inside the
ultraviolet-divergent bare bubble.

The following alternative is not a proved source bridge.
The direct calculation in Appendix~\ref{app:finite-ghost} does not
require this full two-line identity. First identify the ordered-scale ghost graph with the ordinary
two-line integral retaining both factors $f_\rho$, with the same
dimensional continuation and subtraction. At fixed $p\ne0$, its
difference from the uncut integral should be
\begin{equation}\label{eq:ghost-cutoff-difference}
 B_\rho^{\rm ren}(p)-B_\infty^{\MS}(p)
 = C^{ab}\int\frac{d^4k}{(2\pi)^4}
   [1-f_\rho(k)f_\rho(k+p)]
       \frac{k_\mu(k+p)_\nu}{k^2(k+p)^2}.
\end{equation}
This difference is ultraviolet convergent. With
$h_\rho(q)=(1+\rho^2q^2)e^{-\rho^2q^2}$, its absolute integrand is
bounded by
\[
 \frac{h_{\rho_0}(k)+h_{\rho_0}(k+p)}{|k|\,|k+p|},
 \qquad \rho\ge\rho_0>0,
\]
apart from the fixed color and measure factors. This majorant is
integrable in four dimensions: the singularities at $k=0,-p$ are
locally integrable and the tails decay exponentially. Pointwise
convergence and domination would therefore remove the cutoff.
The outstanding formal work is to prove the integral identity with
the submitted continuation, justify these integral operations, and
connect the uncut integral to our Gamma expression. A pointwise
propagator limit or a correct pole alone does not establish that identity.
This is an infrared cutoff in a perturbative graph, not a claim about
uniqueness of the nonperturbative large-volume theory. The completed
test of this individual graph does not replace a test of a complete
physical correlation. \src{S}

There is already a precise negative control on the independent formula.
Freezing $1/(1-\varepsilon)$ to one before multiplying by its pole
would incorrectly replace $8/3$ by $5/3$ in the first tensor channel.
The Lean check proves that the lost finite term is
$\delta_{\mu\nu}p^2/(192\pi^2)$ before the color factor, and that it is
nonzero when that tensor component is one. This is a proved failure of
the altered formula, not yet a mutation run against the submission.

For multiple insertions, \code{gauge_cluster} proves the connected
recursion, and \code{screened} identifies its solution at fixed
infrared length. After regulator removal,
\code{Graphical.identification} identifies \code{fullGraph} with the
full moments of the same \code{MS_Data}. Finally
\code{calibrated_gates} identifies the graph-derived calibrated table
with both \code{Reg.State.wall} and \code{FrameConj.Gates.Thaw}, the
table used on the Wilson side of the asymptotic comparison.
These equalities connect the independent reference calculation to the
matching estimates; they do not define the actual physical correlations
to be graph sums. The latter still come from the lattice and continuum
limits described in Sections~\ref{sec:statement}--\ref{sec:objects}.

\subsection{What this audit establishes and what remains}

The immediate conclusion is stronger than our initial reading:
the submission already contains the all-order graph/flow uniqueness
bridge we wanted to look for. It is therefore inaccurate to describe
the matching result as agreement with unexplained textbook formulas
that were simply entered as definitions. At the same time, internal
agreement does not complete the independent identification of those
definitions with conventional Yang--Mills perturbation theory.

The most useful remaining checks are the extension of the local
ghost/source dictionary to renormalized composite insertions,
the coupling and color conventions, and the identification of the
sector subtraction and evanescent prescription with the intended
$\MS$ scheme. Independent coefficients, especially one outside the
normalization menu, would give a distinct check on this chain.
The signed cubic diagnostic remains valuable, but its role as a
coupling calibration must be kept separate from an additional
prediction. Tests I1--I7 below retain that distinction.

There is a small documentation issue to resolve. Comments in the
graph modules refer to \code{Coverage}, \code{Centrifuge},
\code{ConnectCost}, \code{Spinneys}, and \code{SectorsAtCut}, including
an \code{at_absolute_cut} theorem. We did not locate the corresponding
graph modules or that declaration in the supplied Lean tree (an
unrelated \code{RG02/Flat/Coverage.lean} does exist). We therefore do not cite
the advertised enumeration or unsorting arguments as established
evidence. The checked recursion and uniqueness theorems above are
present; missing comment references do not, by themselves, show a
gap in their proofs.

\code{DiagramReferenceAudit.lean} records twelve submitted theorem
types and their axiom dependencies, together with the two free-line
checks described above. Its build record distinguishes this local
audit against existing compiled submission modules from a new clean
rebuild or Comparator run. It is supporting evidence for the source
reading, not a final independent matching acceptance theorem.
The audit compiled successfully with Lean 4.32.0 and no warnings;
all fourteen axiom reports list only \code{propext},
\code{Classical.choice}, and \code{Quot.sound}.

\code{GhostConventionCheck.lean} adds five proved identities: the two
ghost-line orientations, the cubic source vertex, its equality to the
vertex of the full classical log density, and its scalar-antighost
form. \code{CountertermAudit.lean} records the types and axiom
dependencies of eight submitted pole and finite-part theorems, including
those behind equations~\eqref{eq:actual-pole-complement} and
\eqref{eq:premature-pole}. This second file is an inspection of existing
theorems, not an independent computation of the loop amplitude.
Both files compiled without warnings and reported only the same three
standard axioms. Their source hashes, commands and logs are saved with
the matching verification records.

The subsequent \code{GhostLoopBenchmark.lean} contains eighteen theorem
reports, also with standard axioms only. Its finite reference calculation,
pole comparison, finite-reading bridge, scale-dependence check and
reference-scale reduction pass. It also proves bounds for the heat dressing
and the pointwise infrared limit of the actual integrated line.
The original evidence record keeps that stage's helper proofs distinct
from the then-unproved \code{SubmittedFinitePartTest}.
\code{GhostLoopFiniteEvaluation.lean} adds the explicit radial and
finite parameter calculations described in Appendix~\ref{app:finite-ghost}.
The subsequent modules, ending in \code{GhostLoopSolution.lean}, prove
the full test and connect these calculations to the actual source
graph. The new records \code{evidence/matching/ghost-finite-build.json}
and \code{evidence/matching/ghost-finite-comparator/} certify this
completed comparison; the earlier helper receipts remain unchanged.

\section{Fixing the fields and measuring the coupling}
\label{sec:normalization}

An unrestricted change of field coordinates could make a matching
statement much less informative. The November statement instead fixes
a position-space prescription from the free sectors and gauge group,
then requires the actual physical measurements to obey that prescription.

\subsection{Finite normalization equations for the fields}

For each cutoff $d$, a \code{Screen} is a finite list of separated smooth
compactly supported tests and field tuples, together with a real-linear
way to combine their complex measured values into an endomorphism of
$\mathcal V_{\le d}$. The allowed infinitesimal field changes preserve
the unit and filtration, reality, and signed coordinate permutations.
The free-theory Jacobian of the screen reads every such allowed change
faithfully: applying the screen to the variation of the free moments
returns that change itself. \src{N}

The raw position frame, denoted here by $J_r^{\mathrm{pos}}$, solves
these finite equations so that its screen values equal the free values.
It lies on the branch whose separated test moments tend to the free
Maxwell moments. The statement requires an invertible linearized system
and local uniqueness against permitted changes tending to zero. It
also requires consistent cutoff inclusions and a physical unit in
omitted insertion slots.

This is a local normalization prescription near the ultraviolet
reference. It does not claim global uniqueness of a field coordinate
system at arbitrary physical scale. The menu is selected
noncomputably from admissible menus; the submission does not present
a short numerical list of profiles ready for a laboratory calculation.

\subsection{The dimension-six cubic diagnostic}

The coupling prescription uses two particular polynomial fields.
In an oriented orthonormal frame set
\[
 F_i^{\pm}=\frac{F_{0,i+1}\pm F_{j+1,k+1}}{\sqrt2},
 \qquad(i,j,k)\ \hbox{cyclic}.
\]
Writing $N_{\rm mat}=g.n+1$ for the matrix representation size, the
specified cubic and quadratic polynomials are
\begin{align}
 \cO_6&=\sum_{i,j,k}\epsilon_{ijk}
       \frac{-\Re\operatorname{Tr}\bigl(F_i^+[F_j^+,F_k^+]\bigr)}{N_{\rm mat}},
       \label{eq:dim6}\\
 \cO_4^-&=\frac12\sum_i\langle F_i^-,F_i^-\rangle.
       \label{eq:dim4}
\end{align}
The superscript $-$ here denotes anti-self-duality. It is not a label
for charge-conjugation parity. The code's triangle field named
\code{T} is $\cO_6$, \emph{not} the stress tensor. \src{C}

The three slots are $(\cO_6,\cO_4^-,\cO_4^-)$. Their connected third
moment is formed by the usual subtraction of all proper faces,
including the $2\langle1\rangle\langle2\rangle\langle3\rangle$ term.
Only then is it integrated against one fixed separated profile $\phi$.
Omitted slots carry the unit, so the different faces use the same
geometry and prescription. The physical coupling is exactly
\begin{equation}\label{eq:coupling}
 g_X(r)=\frac{1}{c_T}\Re\int \phi(x)
 \bigl\langle J_r^{\mathrm{pos}}\cO_6(x_1r),
 J_r^{\mathrm{pos}}\cO_4^-(x_2r),
 J_r^{\mathrm{pos}}\cO_4^-(x_3r)\bigr\rangle_{\!c,\mathrm{mixed}}\,dx.
\end{equation}
Here each frame is understood at the appropriate degree; the public
definition includes all three labels in cutoff six.

The real normalization $c_T$ is computed from \code{bareTriangle}.
That expression includes both the nonlinear commutator part of the
inserted curvature and the integrated three-gluon vertex. It contains
the squared structure-constant sum, the chiral contractions, and a
factor six for assignments to the cubic insertion. Derivatives at
the vertex are handled by splitting off cutoff balls and integrating
by parts, including derivatives of the cutoff. This avoids treating
singular differentiated kernels as ordinary absolutely integrable
products without justification. Admissibility requires $c_T<0$.
\src{C}

The raw first connected graph coefficient must integrate to this
\emph{signed} value for every admissible triangle. This tests more than
nonvanishing or the square of a coefficient. The measured coupling is
eventually positive and satisfies
\begin{equation}\label{eq:running}
 g_X(r)\sqrt{2b_0\log(r^{-1})}\longrightarrow1,
 \qquad b_0=\frac{11C_A}{48\pi^2}.
\end{equation}
The Casimir here is defined in the matrix trace normalization
$\langle X,Y\rangle=-\Re\operatorname{Tr}(XY)/N_{\rm mat}$.
For an orthonormal basis $E_a$ the code sets
\[
 C_A=\frac{1}{\dim\mathfrak g}
           \sum_{a,b}\|[E_a,E_b]\|^2.
\]
Numerical comparison with a conventional generator normalization must
include this choice. Equation~\eqref{eq:running} alone is a leading
running law, not a statement of every coefficient of the beta function.
\src{F,N}

\subsection{A single prescribed conversion from \texorpdfstring{$\MS$}{MS-bar}}

The graph reference is converted to the position prescription by a
formal coupling substitution $g_{\MS}=f(g_X)$, with $f(0)=0$ and
$f'(0)=1$, and compatible field series $Z_d(g_X)$ with leading term
the identity. The conversion is determined by the finite normalization
equations and formal inversion of the triangle measurement.
Both orders of the scalar inverse identity and both matrix inverse
identities are required. \src{F,N}

The raw position frame can mix lower engineering weights. A prescribed
strictly lowering correction removes that mixing from the formal
conversion, leaving its equal-weight part. The physical matching frame
is $J_r=J_r^{\mathrm{pos}}A_d(g_X(r))$, where $A_d$ is a fixed smooth
realization of the correction series. The public definition gives this
realization using scaled smooth cutoffs and an explicit sum; it is not
an assertion that the formal series converges analytically.

For correlations, the converted coefficient is the coefficient of
$g_X^k$ after applying $Z$ on every input and substituting $f$ into the
raw graph expansion. For OPEs, one additionally applies $Z^{-1}$ on the
\emph{output} field. This output transformation is necessary for a
field-valued OPE comparison.

The formal conversion is unique for the fixed reference and menu.
The public \code{Reference} is a subsingleton once these requirements
are imposed, including the fixed smooth realization. That is uniqueness
of the perturbative comparison data; it is not uniqueness of the
nonperturbative theory.

\section{The correlation statement in ordinary notation}
\label{sec:correlations}

Let $K\Subset\Conf_n(\R^4)$ be a compact set disjoint from all pairwise
coincidence diagonals. At every truncation order $p$, derivative order
$j$, and fixed collection of degree cutoffs, there is a finite constant
$A$ and a sufficiently small radius such that
\begin{equation}\label{eq:corr-bound}
 \sup_{\substack{x\in K,\ o\in\R^4,\ \sigma\\0\le\ell\le j}}
 \left\|D_x^\ell\left[
 S_{\sigma,r}(v;x,o)-\sum_{k=0}^{p}g_X(r)^k C^X_{n,k}(v;x)
 \right]\right\|
 \le A\,\prod_i\|v_i\|_{\mathrm{coeff}}\,|g_X(r)|^{p+1}.
\end{equation}
The coefficient norm is the finite sum of the absolute values of the
polynomial coefficients. It vanishes only for the zero polynomial and
does not change merely because the same field is presented at a larger
cutoff. The displayed estimate summarizes the stronger
\code{shapeDerivatives} clause. In particular its eventual radius and
constant work uniformly over phases, centers, and the input fields
after the explicit norm factors are included. \src{A}

These are derivatives with respect to the \emph{shape coordinates}
$x$, holding $r$ and the frame fixed. They are not derivatives of
$J_r$ or $g_X(r)$ with respect to $r$. Ordinary physical-coordinate
derivatives introduce the corresponding powers of $r^{-1}$.

The public target also requires smeared expansions: on each fixed
separated compact, the remainder divided by $g_X(r)^p$ tends to zero
uniformly over an appropriate bounded family of smooth compact tests.
Integrability is part of the requirement. The submitted proof obtains
these from the stronger compact estimates and can supply test order
zero at this step. No conclusion about a contact term at coincident
points follows from these separated tests.

Finally, \code{phaseAgreement} gives the additional bound
\begin{equation}\label{eq:phase}
 \|D_x^j(S_{\sigma,r}-S_{\tau,r})(v;x,o)\|
 \le r^\gamma\prod_i\|v_i\|_{\mathrm{coeff}},
 \qquad 0<\gamma<4,
\end{equation}
eventually, uniformly over the indicated phases, centers, fields and
separated compact. Since $g_X(r)$ decays logarithmically,
$r^\gamma=o(g_X(r)^q)$ for every fixed $q$. Thus phase differences
are invisible to every finite coupling order in this simultaneous
short-distance limit. They may still exist at finite physical
separations. Equation~\eqref{eq:phase} should not be read as a statement
about a configuration with arbitrary unscaled distant spectators.

\section{How the correlation proof is assembled}
\label{sec:proof}

\subsection{The construction supplies controlled finite-order expansions}

The native result \code{SignedWilsonExpansion/Result/Orders.lean}
packages finite-order estimates for the actual lattice and continuum
insertions. Its \code{LinePrices} has three parts: a finite-lattice
estimate against a truncated series; comparison of those lattice values
with the selected continuum values; and the resulting continuum
remainder estimate. There are pointwise and smeared forms. The
\code{ConferencePowers.trueLetters} result handles heterogeneous field
degrees in one common cutoff and compares with the independently
nominated subtraction data. \src{P}

The order of limits matters. The native account removes the physical
ultraviolet cutoff with period held fixed, then passes to the chosen
period limit, and finally makes the observation radius small. Thresholds
for the lattice regulators can depend on that radius. This is not a
bound valid for arbitrary independent removal of all regulators.
The construction and its limiting hierarchy are fixed while the
finite perturbative order varies.

This is the analytic input from the renormalization construction.
The present source map locates its use and its precise output. It
does not replace a review of the deeper small-field, large-field and
renormalization estimates that prove this input.

\subsection{The coefficients are identified with an external graph calculation}

Having an expansion into some formal coefficients would not suffice.
The second connection is a coefficient-by-coefficient equality with
the graph reference of Section~\ref{sec:graphs}. The source proves that
the graph calculation and the Wilson-side perturbative calculation,
after the same position calibration, are limits of a common regulated
calculation.

In \code{original_calibrated_gate_of_team}, the last step is uniqueness
of a limit: one convergence theorem gives the continuum graph reading,
another gives the Wilson-side reading, and the regulated expressions
are identified using compatible prescriptions. The theorem
\code{original_filed_calibrated_gates} supplies the compatible choices.
\code{Graphical.CommonGates.calibrated_gates} identifies the explicit
graph sum with these calibrated readings. This is an actual equality
argument, not a definition that calls the coefficients of the Wilson
answer ``Feynman diagrams.'' \src{E}

The subsequent public bridges include mixed-degree words in a common
sector and identify each coefficient after the prescribed coupling and
field substitution. \code{publicMixedPolynomial} identifies the finite
polynomial used in the analytic estimate with the finite sum of the
converted graph coefficients. The proof then substitutes this equality
into the remainder estimate. \src{P}

\subsection{Position derivatives require an additional argument}

A uniform estimate on functions alone does not justify differentiating
an asymptotic expansion. Here the derivative proof uses the physical
field derivative maps: differentiating an insertion raises its field
degree, and the raised insertion can be expressed in the next finite
field sector. The coefficients of this operation have their own
controlled stationary expansions.

\code{WilsonGraphCompactCoordinateJets.prices} inducts on a word of
coordinate derivatives. At each step it combines the expansion of a
raising coefficient with the expansion of the correlation containing
the raised field. The accompanying smoothness and
\code{cropCoordinates} theorems identify the result with derivatives
of the genuinely truncated coefficient functions. Finite-dimensional
coordinate estimates are then assembled into the Fr\'echet-derivative
bound. Translation covariance permits arbitrary physical centers,
and multilinearity extends basis estimates to all input fields with
the norm factors in \eqref{eq:corr-bound}. \src{J}

\subsection{One hierarchy, all its phases, and the public theorem}

The native estimates refer to the actual mixed hierarchy and its own
compression characters. The field identification used by
\code{WilsonGraphPublishedShapeEstimate.onField} holds on a whole
separated neighborhood, so it transports derivatives as well as values.
The stronger phase proximity estimate uses the low-degree field content
and the short-distance OPE: below degree four the only surviving
physical field is the unit, whose normalized expectation is the same
in every phase. Remainders yield positive powers arbitrarily close
to four; the frame bounds control the cost of changing coordinates.
\src{J,H}

Finally \code{graphEstimates} transfers these results to the public
coefficient array, and \code{correlationInFrame} supplies the
\code{CorrelationExpansion} in \code{MatchedRefinement}. The assembly
also retains the native lattice coordinates, fixed-period kernel limits,
and the same selected hierarchy. An additional compatibility condition
says that, \emph{if} a requested volume sequence already has full tested
limits, the chosen hierarchy represents those limits. This conditional
clause does not prove convergence of every requested volume sequence.
\src{T,V,J}

\section{The OPE: characterization, uniqueness, and matching}
\label{sec:ope}

\subsection{What determines the reference OPE}

The independent reference OPE is characterized using the graph
\emph{correlations} and arbitrary exterior fields. For homogeneous
inputs of total engineering degree $D$, shrink their positions by
$\lambda=e^{-t}$ while exterior insertions $b$ at positions $y$ remain
fixed. The defining statement has the form
\begin{align}\label{eq:exterior}
 C^{\MS}_{n+m,k}(a,b;e^{-t}x,y)
 ={}&\sum_{w=0}^{h}\sum_{q=0}^{k}
 e^{(D-w)t}
 C^{\MS}_{m+1,k-q}\bigl(p_{w,q}(t,x),b;0,y\bigr)
 +R_{h,k}(t,x,y),\\[-2pt]
 |R_{h,k}(t,x,y)|\le{}&c(1+t)^L e^{(D-h-1)t},
 \qquad t\longrightarrow\infty.\nonumber
\end{align}
The bound is uniform on separated compact sets for the inner shape and
for the exterior configuration, the latter separated from the fusion
center. Each $p_{w,q}$ is homogeneous of output weight $w$ and
polynomial in $t$, coordinate by coordinate, with a finite degree bound
at each order. The public reference output through degree $h$ is
$\sum_{w\le h}p_{w,k}(0,x)$. Extension to nonhomogeneous inputs is
multilinear and independent of the cutoff used to present a polynomial.
\src{O}

The predicate \code{Ordinary.Determined} asserts existence and
uniqueness of this entire OPE array. The proof does more than choose
one convenient expansion: \code{GraphOnlyPrice.original} compares two
candidate exterior expansions, turns their difference into a residual
in the native exterior system, and uses that system's uniqueness
result. \code{GraphOrdinary.original} then packages the unique full
reader. Thus positive-weight output channels are tested by exterior
correlations, not only by their vacuum expectation. Reviewing the
exterior system's separation and uniqueness argument remains one of
the deeper proof-review tasks. \src{O}

\subsection{What physical OPE is compared}

On the constructed hierarchy, the actual OPE has channels of each
degree with physical channel fields and numerical coefficient
functions. The public \code{measuredOPE} sums every channel through
the output cutoff $h$, includes its field in the common physical
space, and applies $(J_r^{(h)})^{-1}$ to the result. Call this
$\mathcal C_r^{\le h}(a;x,o)\in\mathcal V_{\le h}$.

The predicate \code{HasChannelBasis} requires these included channel
fields to form a basis of the \emph{whole} output space through $h$.
The final matching estimate is
\begin{equation}\label{eq:ope-bound}
 \left\|\mathcal C_r^{\le h}(a;x,o)
       -\sum_{k=0}^{p}g_X(r)^k C^{X,\le h}_{\mathrm{OPE},k}(a;x)
 \right\|_{\mathrm{coeff}}
 \le A\prod_i\|a_i\|_{\mathrm{coeff}}|g_X(r)|^{p+1}
\end{equation}
on each separated compact. There are corresponding smeared estimates
against every linear functional of the output and shape-derivative
estimates uniform over the dual unit ball. The derivative version is
also uniform in the center. The OPE reader itself has no phase index;
the correlation hierarchy and the separate exterior estimates explain
how that reader is used in phases. \src{A}

The two truncations are different. Increasing $p$ improves the
perturbative expansion at fixed finite output cutoff $h$. Increasing
$h$ retains more physical fields in an OPE expansion whose geometric
remainder is controlled separately. Neither statement asserts that
the perturbative series converges as $p\to\infty$.

\subsection{Why correlation matching alone is not the whole proof}

To recover an OPE output one must distinguish fields by how they pair
with exterior insertions. The proof therefore controls two kinds of
correlation: the inner cluster paired with exterior fields, and the
correlation with a single output field at the fusion center paired
with the same exterior fields. The latter supplies the finite-dimensional
comparison needed to read the former as an output field. Errors in
this comparison must be controlled, not merely its formal invertibility.

The modules culminating in \code{AppointedFlashNumerator.allJets}
establish this comparison in a window with two relative length scales.
First they prove estimates on fixed basis words. A finite sum extends
these to bounded sets of input and exterior fields, with a common
polynomial logarithmic loss. The source combines the cluster estimates
with separate centered-field estimates in
\code{measuredAtPositiveCrops}. \code{AppointedRoundedReader.value}
then bounds the actual field-valued OPE remainder; further derivative
arguments give all shape derivatives. \src{Q}

At the final interface,
\code{WilsonFlashScoredOrdinaryEstimates.estimates} converts the
uniform bounds against contractive output functionals into a norm
bound by choosing a functional attaining the norm of the residual.
\code{ordinaryInFrame} identifies the reference reader with the unique
public OPE array and supplies \code{OPEExpansion}. This last norm step
is why a scalar vacuum comparison cannot substitute for the stated
full output comparison. \src{Q,T}

\section{The stress tensor has its own matching requirement}

The reference contains a unique degree-four field series
$P_{ab}(g_X)=\sum_{k\ge0}g_X^k P_{ab,k}$. Its leading term is the free
metric stress tensor. The defining conditions require symmetry,
signed-axis covariance, vanishing one-point functions, conservation,
and a translation Ward identity against general field words, with
the prescribed field and coupling conversion applied to all legs.
The zero one-point condition fixes the vacuum ambiguity. \src{W}

The translation identity uses a cutoff window and stress flux. It fixes
the action and sign of the stress insertion on other operators. For
every finite order~$p$ and linear output functional~$\ell$, matching
identifies the \emph{actual} physical stress tensor in the same frame:
\begin{equation}\label{eq:stress}
 \frac{\ell\bigl((J_r^{(4)})^{-1}(r^4T_{ab})\bigr)
       -\sum_{k=0}^{p}g_X(r)^k\ell(P_{ab,k})}
      {g_X(r)^p}\longrightarrow0.
\end{equation}
This is the public \code{stressExpansion} clause. Its displayed
little-$o$ form should not be silently replaced by the stronger
uniform $O(g_X^{p+1})$ derivative bounds stated for general correlations.
It also should not be confused with the cubic field used in the
coupling triangle. \src{A,W}

\section{What this adds, and what still needs human validation}
\label{sec:validation}

The November matching target materially strengthens the evidence for
the intended Yang--Mills theory. A free limit tests order zero. A
leading stress-trace two-point coefficient tests one specific piece
of short-distance information. The present target fixes a nonabelian
cubic diagnostic with its sign, all coefficients of all separated
correlations in the stated local field system, and all finite-degree
OPE outputs, after one prescribed scheme conversion.

The correct assessment is nevertheless not that every definition has
already been independently endorsed. The following checks organize the
next review. ``Traced'' means that the definitions and proof connections
were located and inspected for this document; it is not a substitute
for a second mathematical derivation.

The selected definitions are now collected in
Appendix~\ref{sec:glossary}. Section~\ref{sec:inventory} separates
the supplied theorem obligations from our proposed independent
tests in Section~\ref{sec:independent}. A definition being listed
or a test being specified does not mark that check as completed.

The most useful first independent calculation is the cubic diagnostic
already built into the public prescription. It directly tests a
nonabelian interaction, has an explicit signed formula, and addresses
our earlier interest in a dimension-six curvature operator. A second
calculation should test a positive-weight OPE output or a low-order
correlator not used to define the scheme. Agreement for a defining
normalization observable has a different evidential role from agreement
for an additional prediction.

There are several limits to the conclusion. The theorem gives an
asymptotic expansion at each fixed order, with constants and radius
thresholds allowed to depend on that order. It gives no uniform
large-order estimate, Borel summability, convergent perturbation
series, or effective numerical error algorithm. Equal expansions
permit differences smaller than every power of the coupling; the
phase bound gives a specific short-distance instance of this
possibility. Infrared phase uniqueness, a mass gap, and a unique
large-volume theory are not consequences of the matching result.

The separated-position statements also leave contact terms outside
their direct scope. Graphs may require subtraction at internal
coincidences, and the stress Ward statement controls its specified
flux tests, but these facts do not turn the correlation/OPE matching
clauses into equality of unrestricted contact extensions. Finally,
this document does not claim a separate formal no-go theorem excluding
every possible realization as composites of Gaussian fields. Any
such conclusion should be tied to a precisely defined alternative
and a coefficient or structural condition it fails.








% BEGIN REVIEW_TESTS

\section{Inventory of the tests supplied by OpenAI}
\label{sec:inventory}

The submission supplies two public Comparator configurations:
\code{ComparatorChallenges/Existence.json} checks
\code{YangMills.jointExistence}, and
\code{ComparatorChallenges/Matching.json} checks
\code{YangMills.perturbativeMatching}. Both allow only
\code{propext}, \code{Quot.sound} and \code{Classical.choice};
both set \code{enable_nanoda} to \code{false}.
The matching witness also implies existence for that same run.
The recorded results of checking the Matching target are in
Section~\ref{sec:receipts}.

The list below decomposes the Matching target into fifteen review
groups. These are obligations proved together by the theorem,
\emph{not fifteen separately executed numerical tests}. A passed
Comparator check covers the formal definitions of these obligations;
it does not independently endorse their physical interpretation.
The adjacent \code{openai-test-inventory.json} indexes the individual
fields of the selected public structures and preserves the exact
two configurations. The labels G01--G24 refer to the glossary in
Appendix~\ref{sec:glossary}.

\subsection*{O01. The specified microscopic theory}
\code{MatchedRun.lattice}, \code{trajectory} and \code{admissible}
attach the witness to admissible Wilson data and native insertion coordinates.
Read G04, G07 and G08 together; the lattice carrier alone is not the full test.

\subsection*{O02. Full fixed-period ultraviolet convergence}
\code{meshLimit} and \code{fixedPeriodLimits} identify one smooth
period kernel with the full mesh sequence. All finite difference
stencils, compact interpolation and suitable smearings are included (G09).

\subsection*{O03. The same hierarchy after the volume limit}
\code{MatchedRefinement.continuum} includes the lattice-limit,
field-space, phase, reconstruction, stress and ultraviolet laws.
\code{testLimits} represents any test table already convergent
on the requested sequence. It does not assert unconditional full-volume
convergence (G05--G07, G09).

\subsection*{O04. Actual physical OPE}
\code{singleCluster} and \code{fusionTrees}, together with the
continuum OPE and coefficient laws, give the physical single-cluster
and nested expansions. The matching clauses below concern those
same physical coefficients, not an auxiliary fitted array.

\subsection*{O05. A reference fixed by the gauge group}
\code{Reference.opeDetermined} and \code{groupMenu_eq}
fix the ordinary OPE and normalization menu from the group-only
graph reference. Reference uniqueness is proved (G10--G11, G21).

\subsection*{O06. The signed cubic coefficient}
\code{Calibration.signedSlope} requires
\code{Triangle.RawSlopeAgrees} for every admissible triangle.
This compares an explicit bare formula with the first connected
graph coefficient, including all proper-face Leibniz terms (G13--G14).

\subsection*{O07. Formal normalization and inverse transformations}
\code{Calibration.conversionValid} includes the raw normalization
equations, vanishing free offset, nonzero signed slope, both inverse
coupling identities, both inverse field series and the factorization
of the strictly lowering correction. It determines the conversion
coefficient by coefficient (G12, G15).

\subsection*{O08. Numerical normalization beyond formal jets}
\code{adjustmentValue} identifies the numerical correction with
\code{smoothAdjustment}. The \code{NumericalAdjustment} laws impose
its smoothness, asymptotic expansion, local invertibility, reality,
unit and filtration properties (G16).

\subsection*{O09. Calibration on actual continuum observables}
\code{FrameCalibration} requires physical position normalization,
local uniqueness, reality, unit and nesting, triangle integrability,
the exact measured-coupling equation, and the relation between the
position and measured frames (G17--G18).

\subsection*{O10. One coherent measured field system}
\code{Measurement} requires compatible fields across cutoffs,
signed-axis covariance, unit insertion in every phase, compatible
field inclusions and correlation multilinearity (G18).

\subsection*{O11. Positive coupling and the one-loop law}
\code{Measurement.coupling_pos} and \code{oneLoop} require
$g_X(r)>0$ eventually and
$g_X(r)\sqrt{2b_0\log(1/r)}\to1$.
\code{oneLoopBeta} computes $b_0$ from trace-normalized
commutators. This particular clause is a leading-log statement,
not a differential beta equation.

\subsection*{O12. Free ultraviolet limit at all radii}
\code{scaling}, \code{sampledFrame}, \code{fixedFields},
\code{maxwellLimit} and \code{frameBounds} connect the sampled
Maxwell limit to the same measured fields at all small radii,
with the additional fixed-field and frame control (G24).

\subsection*{O13. All-order correlation matching and phase proximity}
The first four fields of \code{CorrelationExpansion} require
the smeared, compact and derivative expansions and the
$r^\gamma$ phase-difference estimate. They quantify over
arbitrary local labels and the stated phases (G19--G20).
The reference is common; neither the constants nor the expansion
are required to converge as perturbative order grows.

\subsection*{O14. All-order matching of every finite OPE output}
\code{channelBasis} and the three \code{OPEExpansion} clauses
require a full field-space basis and the smeared, norm and
derivative comparisons after inverse output conversion (G22).
The physical output has no separately chosen phase argument.

\subsection*{O15. Stress Ward normalization and physical identification}
\code{Reference.stressUnique} fixes a unique series satisfying
all six \code{WardStress.Determined} laws.
\code{CorrelationExpansion.stressExpansion} identifies the
actual physical stress with that series (G20, G23).

\section{Independent tests: precise first specifications}
\label{sec:independent}

Our tests should expose a physical observable on the left and an
independently calculated quantity on the right. Repackaging
\code{CorrelationExpansion} or evaluating the same graph function
on both sides would be a useful interface check, but would not
validate a coefficient. The strategy is to derive a few low-order
numbers or functions from explicit propagators and interactions,
translate their conventions to the public fields, and then use
the submitted matching theorem to identify the actual continuum
prediction.

The accompanying \code{IndependentMatchingTests.lean} starts this
work. It defines independent kernel and coefficient specifications,
an entire-output OPE remainder specification, and a finite-step
running-coupling test. It also exposes the actual submitted
correlator through a small adapter. The finite Pauli-matrix
identities described below are proved. The adapter equality is
only a definitional sanity check. No combined independent acceptance
theorem is yet proved, and no new Comparator run on these test
specifications has been performed. The separate full finite ghost-loop
benchmark has passed in Lean and Comparator, as described in
Section~3.5 and Appendix~\ref{app:finite-ghost}.

\subsection{I1. Fix the free kernel and color normalization first}

For $x\ne0$ define independently
\[
 G(x)=\frac{1}{4\pi^2|x|^2},\qquad
 H_{ab}(x)=\frac{1}{4\pi^2}
 \left(\frac{8x_ax_b}{|x|^6}-\frac{2\delta_{ab}}{|x|^4}\right).
\]
The curvature covariance to check is
\[
 K_{ab,cd}(x)=\delta_{bc}H_{ad}-\delta_{bd}H_{ac}
             -\delta_{ac}H_{bd}+\delta_{ad}H_{bc}.
\]
This includes the minus sign from differentiating the second
insertion. \code{expectedHessian}, \code{expectedCurvature} and
\code{freeKernelTest} state the formulas without calling the
submitted covariance. The next analytic lemma is that the
submitted double curl equals $K$ off the origin, followed by
the comparison of its shell-coordinate derivatives, factorials
and Wick multiplicities. No claim about a contact extension
at $x=0$ is implicit here.

There is already a small completed algebra check. Take the
anti-Hermitian Pauli basis $E_a=-i\sigma_a$, in the displayed
order in the Lean file. The independent calculation proves
\[
 -\tfrac12\operatorname{Re}\operatorname{Tr}(E_aE_b)=\delta_{ab},
 \qquad [E_a,E_b]=2\epsilon_{abc}E_c.
\]
Consequently $\sum_{abc}(2\epsilon_{abc})^2=24$ and
$C_A=24/3=8$ in this color metric.
The value is deliberately expressed in the public normalized-trace
convention, not in the frequently used convention
$\operatorname{Tr}(T_aT_b)=\delta_{ab}/2$.
The remaining bridge must identify the actual submitted color
frame with this convention, up to an orthogonal basis change,
before substituting the result into \code{oneLoopBeta}.

\textbf{Status:} finite matrix and color identities checked in Lean.
The separate \code{DiagramReferenceAudit.lean} now also checks the
momentum-space gauge line and its cutoff factor
(equation~\ref{eq:gauge-line-audit}). Position-space kernel
differentiation and the public-frame identification remain open.
This is a normalization baseline, not yet an independent interaction test.

The ghost-loop benchmark in Section~3.5 is a completed additional
check: the full submitted finite graph agrees with the independently
specified expression, including both finite constants. Its Lean proof
and separate Comparator run passed. Appendix~\ref{app:finite-ghost}
explains how the source integral and its infrared limit are evaluated.
This graph test remains separate from the physical correlation tests
I2--I4.

\subsection{I2. Calculate the signed dimension-six cubic diagnostic}

Use the public polynomials, not just their names:
\[
 \cO_6=\sum_{ijk}\epsilon_{ijk}
 \frac{-\operatorname{Re}\operatorname{Tr}
       (F_i^+[F_j^+,F_k^+])}{N},
 \qquad
 \cO_4^-=\frac12\sum_i\langle F_i^-,F_i^-\rangle .
\]
The minus sign in $\cO_4^-$ means anti-self-duality.
For the convention $\exp(-S_0-gS_3-g^2S_4)$, the first coefficient
of the connected three-point function has two types of terms:
\[
 \left\langle
   \left.\partial_g
     \bigl(\cO_6(g)\cO_4^-(g)\cO_4^-(g)\bigr)\right|_{g=0}
 \right\rangle_{0,c}
 -
 \left\langle\cO_6(0)\cO_4^-(0)\cO_4^-(0)\,S_3\right\rangle_{0,c}.
\]
Here the notation means the coefficient obtained by differentiating
the normalized connected expectation, including its proper-face
terms. The first term matters because the curvature itself
contains the commutator of gauge potentials. A calculation of
only the three-gluon vertex would not check the submitted formula.
The independent calculation must also track the cutoff partition,
integration by parts, and any subtraction required by its chosen
intermediate representation.

First compare this explicit Wick calculation with
\code{bareTriangle} and with the raw first graph coefficient.
Then perform the independently understood field and coupling
conversion to compare a physical prediction. The generic
\code{firstCoefficient} specification asks for
\[
 \lim_{r\downarrow0}\frac{A(r)-A_0}{g_X(r)}=A_1^{\rm independent},
 \qquad g_X(r)>0,\quad g_X(r)\longrightarrow0.
\]
The elected triangle itself defines $g_X$ exactly. Its normalized
first slope is therefore not an additional physical prediction.
Use another admissible profile or another separated observable
after checking the raw coefficient and conversion. OpenAI's
every-admissible-triangle condition is useful here.

\textbf{Status:} the coefficient-extraction and connected-three
specifications are written in Lean. The independent spacetime
calculation and its identification with the source are not yet proved.

\subsection{I3. Test an order-two prediction outside the calibration menu}

Choose four separated gauge-invariant quadratic curvature fields,
in a fixed homogeneous normalization, and compare their connected
correlation through order $g_X^2$. The intended requirement is
\[
 \frac{A(r)-A_0-g_X(r)A_1}{g_X(r)^2}
       \longrightarrow A_2^{\rm independent}.
\]
\code{secondCoefficient} states this limit; the more general
\code{finiteOrderRemainder} also records a compact uniform
remainder against an independently supplied coefficient array.
The actual field labels and a convenient separated configuration
must still be chosen.

The order-two calculation includes the quartic vertex, the
$S_3^2/2$ contribution, nonlinear curvature terms, ghost terms
where they occur, normalization of the expectation, and the
prescribed counterterms. It is not just a quartic tree amplitude.
Moreover a connected four-point function of quadratic Gaussian
composites can already be nonzero at order zero. The test must
subtract that known Wick contribution, rather than equating
nonzero connected correlations with a non-Gaussian elementary field.

\textbf{Status:} a precise coefficient-extraction specification is
formalized; the chosen observable, calculation and physical bridge
remain to be supplied. This is intended to test a prediction
not used to set the coupling.

\subsection{I4. Check a nonidentity OPE channel as a field}

Start with two quadratic curvature scalars and the full output
space through degree four. Compute the free coefficient first,
then the first nonvanishing interaction correction. Fix exterior
probe words whose free correlation matrix detects that finite
output space. Its invertibility must be established; otherwise
a vanishing reading need not mean a vanishing field.

In the measured frame the proposed test has the form
\[
 \sup_{x\in K}
 \left\|\cO^{\rm actual}_{\le4}(r,x)
       -\sum_{k=0}^p g_X(r)^k
           \cO^{\rm independent}_{\le4,k}(x)\right\|
       \le C_K |g_X(r)|^{p+1}.
\]
The generic Lean predicate is \code{finiteOutputRemainder}.
The expected output includes positive-weight fields, not just
the identity coefficient. The comparison must apply the inverse
field conversion on the output as well as the conversions on
the two inputs. The submission's \code{HasChannelBasis} and
exterior uniqueness are the natural bridge to its actual OPE.

\textbf{Status:} the field-valued specification is formalized;
concrete probes, their detecting matrix, the chosen norm bridge
and independent coefficients remain open.

\subsection{I5. Check stress action, including its sign}

Choose a nonconstant separated field word and a smooth window
isolating one insertion. Derive its translation identity from the
chosen Euclidean action and stress convention. In schematic form,
\[
 \sum_a\int (\partial_a\chi)(z)
       \langle T_{ab}(z)\cO(x)\cdots\rangle\,dz
   =\sigma_{\rm conv}\,\partial_{x_b}
       \langle\cO(x)\cdots\rangle .
\]
The sign $\sigma_{\rm conv}$ is to be fixed in advance by that
derivation, not fitted to the submitted answer. The actual test
will smear the other variables and use the source's precise
window clearance. Comparing it with the two minus signs and
the derivative of the test function in
\code{WardStress.Determined.translation} will check the
integration-by-parts convention as well.

Use the same stress field that appears in physical
\code{stressExpansion}. A trace two-point function or a
conservation law by itself does not determine the translation
generator's sign and normalization on a general operator.

\textbf{Status:} the source obligation and a concrete test plan are
identified; this independent window identity has not yet been
written as a Lean theorem.

\subsection{I6. Test running between two physical scales}

For a fixed ratio $s>1$ of physical lengths, ask for
\[
 \lim_{r\downarrow0}
 \frac{g_X(sr)-g_X(r)}{g_X(r)^3}=b_0\log s.
\]
\code{oneLoopStepScaling} states this requirement independently.
The positive sign is for increasing length; it is opposite to
the energy-scale beta convention. Unlike a mere leading-log
equivalent, this condition resolves a small difference between
nearby scales. The law
$g_X(r)\sqrt{2b_0\log(1/r)}\to1$ alone does not imply it.

The native result \code{ConferencePowers.guideNext} is a promising
place to look for sufficiently precise step control. Its passage
to the measured physical coupling still has to be proved.
This is an additional diagnostic, not a claim that the public
Matching target explicitly promises this stronger formulation.

\textbf{Status:} the finite-step specification is formalized;
the theorem connecting it to the construction is open.

\subsection{I7. Check that the independent tests reject deliberate changes}

After the preceding calculations have concrete outputs, make
small controlled alterations: reverse the cubic sign; change the
quartic coefficient or omit a ghost contribution where it contributes;
rescale the color bracket without the corresponding metric and
coupling change; retain only the identity OPE output; omit
$Z^{-1}$ on its output; or reverse the stress action. Each
alteration should fail a named independent equality with a
visible reason. A consistent change of conventions should pass
after the documented conversion.

An elementary generalized free field can be tested through its
higher cumulants. Composites of Gaussian fields require comparison
with their actual Wick cumulants, not the elementary-field
criterion. We do not yet have a theorem excluding every possible
Gaussian-composite realization.

\textbf{Status:} these negative controls are proposed, not run.
No counterfeit has been shown to pass the November target.

\subsection{Order of work and acceptance status}

Complete I1's bridge to the submitted covariance and color
normalization first. Then calculate I2's cubic coefficient.
That gives a controlled nonabelian test with conventions fixed.
I3 and I4 should provide predictions beyond the defining
normalization observable, while I5 checks the actual stress
action. I6 is useful additional running-coupling information.
Run I7 only after the corresponding expected quantities are fixed.

This is a starting set of independently formulated tests, not a
claim that passing seven selected examples validates every
all-order graph definition. They complement a systematic
definition review and can expose errors that a successful
Comparator check against a misinterpreted definition would not.
The adjacent \code{independent-tests-build.json} records the
small Lean-file check separately from the submission's
Comparator receipts.

% END REVIEW_TESTS

\section{Mechanical verification already completed}
\label{sec:receipts}

The Matching challenge passed the release-pinned Comparator on
10 November at 05:30:13 UTC. It also passed a repeat with the newer
upstream Comparator at 08:18:57 UTC. Both records have exit code zero.
The latter is 00:18:57 Pacific time on 10 November. The axiom audit
reports only \code{propext}, \code{Classical.choice}, and
\code{Quot.sound} for the checked matching theorem. \src{Z}

The release-pinned Comparator revision was
\code{07bc4ea40f2266dcb861820a2ec1fa3244ed307f}; the newer revision was
\code{d03acab154d269c06e60e4de7e4cc85deebff94b}. The submitted proof
used Lean 4.32.0; the repeat used the Lean 4.34.0 checking kernel and
a compatible export of the submitted environment. The runs used
macOS development mode with the \code{fake-landrun.sh} launcher,
not Linux/Landrun isolation or an independent second kernel.

A separate run on 13 November checked our independent full finite
ghost-loop target, using the newer Comparator and the same toolchain
pair. It passed with no definition holes and only the three permitted
standard axioms. Its evidence is kept separately from the supplied
Matching target: the former checks one explicit graph value, whereas
the latter checks the submitted all-order matching statement.
Appendix~\ref{app:finite-ghost} gives the proof map and frozen-target
records.

Comparator checks the submitted proof against a frozen target and its
definitions. Its success is evidence for the formal implication.
Human review must still decide whether the frozen graph definitions,
normalization prescription, and quantified estimates express the
intended physical comparison. The separation between these tasks is
especially important for this larger target.

\appendix

% BEGIN REVIEW_GLOSSARY

\section{Glossary of the key formal definitions}
\label{sec:glossary}

The entries follow the flow from the gauge group and lattice
fields to continuum observables and their perturbative comparison.
Each gives an informal meaning and the exact formal declaration
and source location. Paths are relative to the verified November
source tree recorded in \code{source-map.json}.
This is a glossary of the principal interfaces, not a count of
every definition in the proof's dependency closure.
\code{definition-glossary.json} preserves the same list with
source hashes and declaration lines.

\Needspace{8\baselineskip}
\subsection*{G01. Gauge group, representation and spacetime}
\label{gloss:G01}
\textbf{Formal:} \code{YangMills.GaugeGroup}.\newline
\textbf{Source:} \code{YangMills/Geometry.lean:25}.

The input includes a closed subgroup of a unitary matrix group and its exact
one-parameter Lie algebra. \code{IsSimple} requires real simplicity and
noncommutativity. Matrix trace conventions enter the free color metric and
therefore every coupling coefficient. Spacetime has four axes. The distinguished
\code{timeAxis} is index~3, whereas the chiral formula below uses the oriented
0123 convention with axis~0 singled out. Tracking that change of convention is
part of the sign review; it is not by itself an inconsistency.

\Needspace{8\baselineskip}
\subsection*{G02. Free on-shell curvature coordinates}
\label{gloss:G02}
\textbf{Formal:} \code{YangMills.Maxwell.Coordinates}.\newline
\textbf{Source:} \code{YangMills/FreeFields.lean:30}.

At derivative order $j$, the shell consists of symmetric derivative tensors
of an antisymmetric two-form satisfying the free Bianchi and divergence equations.
A basis supplies polynomial coordinates $(j,a)$; $a$ indexes the chosen shell
basis, not a spacetime index. Each coordinate has engineering weight $j+2$.
The color map obeys $-\operatorname{Re}\operatorname{Tr}(E_AE_B)=N\delta_{AB}$,
where $N=g.n+1$. The covariance differentiates the massless double curl and
includes Taylor factorials and the sign of derivatives on the second insertion.

\Needspace{8\baselineskip}
\subsection*{G03. Maxwell field reference}
\label{gloss:G03}
\textbf{Formal:} \code{YangMills.Maxwell.Reference}.\newline
\textbf{Source:} \code{YangMills/FreeFields.lean:99}.

The degree-$d$ sector consists exactly of the gauge-invariant polynomials
of weight at most $d$ in the preceding coordinates. Moments are fixed by all
complete cross-insertion Wick pairings. Self-pairings at a single insertion
are excluded. The reference also fixes the free stress field and requires
that exterior correlations detect every nonzero sector element. This
\code{Maxwell.Reference} is the free field system; it is distinct from the
perturbative \code{Reference} in G11.

\Needspace{8\baselineskip}
\subsection*{G04. Microscopic Wilson family}
\label{gloss:G04}
\textbf{Formal:} \code{YangMills.WilsonFamily}.\newline
\textbf{Source:} \code{YangMills/WilsonLattice.lean:35}.

A family records physical periods, inverse mesh sizes, the bare coupling,
Haar measure and a token-indexed collection of lattice insertions.
Its finite moments integrate those insertions with the specified Wilson
weight. The matching target additionally requires a family arising from
admissible native lattice data, not an arbitrary assignment of token values.
The ultraviolet index and the physical-period index are separate parameters.

\Needspace{8\baselineskip}
\subsection*{G05. Continuum correlation hierarchy}
\label{gloss:G05}
\textbf{Formal:} \code{YangMills.Hierarchy}.\newline
\textbf{Source:} \code{YangMills/CorrelationHierarchy.lean:24}.

This is the common carrier of continuum fields, correlations, phase
characters, stress tensor and OPE data. A correlation has field labels,
positions and an optional character argument: \code{none} is the mixed
hierarchy, and \code{some character} is a selected phase. The carrier alone
does not prove it is a Wilson limit. The continuum, phase, reconstruction
and matching laws attached to the same hierarchy provide that connection.

\Needspace{8\baselineskip}
\subsection*{G06. Construction and volume refinement}
\label{gloss:G06}
\textbf{Formal:} \code{YangMills.Wilson.Construction}.\newline
\textbf{Source:} \code{YangMills/VolumeLimit.lean:14}.

A construction contains one Wilson family and a refinement operation.
Given a strictly increasing sequence of volume indices, the operation
returns a strictly increasing subsequence and a hierarchy. Existence and
matching predicates supply the laws on those choices. This is selection
from every requested volume sequence, not a claim that the unselected
large-volume sequence converges or that different phases coincide.

\Needspace{8\baselineskip}
\subsection*{G07. The full matching witness}
\label{gloss:G07}
\textbf{Formal:} \code{YangMills.PerturbativeMatching.MatchedRun}.\newline
\textbf{Source:} \code{YangMills/PerturbativeMatching.lean:44}.

One reference, lattice data, period kernels and trajectory serve all
requested volume sequences. Each \code{MatchedRefinement} contains the
continuum laws, measured fields and coupling, calibrated frame, channel
basis, Maxwell limits, and correlation and OPE expansions on its selected
hierarchy. The theorem asks for this entire dependent package.
\code{MatchedRun.realizes} proves that the same run satisfies the existence
requirements.

\Needspace{8\baselineskip}
\subsection*{G08. Native insertion coordinates}
\label{gloss:G08}
\textbf{Formal:} \code{YangMills.PerturbativeMatching.NativeCoordinates}.\newline
\textbf{Source:} \code{YangMills/PerturbativeMatching/NativeLattice.lean:26}.

Raw field coordinates are the coefficient vectors of the original
graded lattice insertions, in their original order. The width equality
identifies the finite coefficient spaces, and the token equality identifies
their actual insertions. \code{AdmissibleLattice} combines this with the
native family's admissibility laws. These requirements matter when
checking that the later physical fields remain attached to the specified
microscopic observables.

\Needspace{8\baselineskip}
\subsection*{G09. Smooth full-sequence ultraviolet limit}
\label{gloss:G09}
\textbf{Formal:} \code{YangMills.PerturbativeMatching.SmoothMeshLimit}.\newline
\textbf{Source:} \code{YangMills/PerturbativeMatching/VolumeLimits.lean:67}.

For each fixed period and token word there is one smooth kernel away
from periodic coincidences. Every finite list of lattice difference
stencils tends to its corresponding derivative, including evaluation
along convergent mesh points, compact uniform interpolation and smearing.
The empty stencil recovers the original moment. Its restriction agrees
with the period table used for the volume limit. The separate
\code{RepresentsTestLimits} clause is conditional on a whole test table
already converging along the requested volume sequence.

\Needspace{8\baselineskip}
\subsection*{G10. Continuum graph coefficient}
\label{gloss:G10}
\textbf{Formal:} \code{YangMills.PerturbativeMatching.graphCorrelation}.\newline
\textbf{Source:} \code{YangMills/PerturbativeMatching/GraphCoefficients.lean:43}.

This function supplies the coefficient of a specified perturbative
order for a word of free polynomial labels at separated positions.
Its inputs are the gauge group and free labels, not a Wilson hierarchy
or phase. It is assembled from the continuum graph rules, pole
subtractions and regulator limits described in Sections~3 and~6.
The human validation task is to identify those rules with the intended
Yang--Mills perturbation theory, including signs and normalization.

\Needspace{8\baselineskip}
\subsection*{G11. Calibrated perturbative reference}
\label{gloss:G11}
\textbf{Formal:} \code{YangMills.PerturbativeMatching.Reference}.\newline
\textbf{Source:} \code{YangMills/PerturbativeMatching/Reference.lean:35}.

This package requires the graph correlations to determine a unique
ordinary OPE, chooses the group-only normalization menu, calibrates a
position scheme and requires a unique Ward-normalized stress series.
Its OPE is read from the exterior characterization, not freely supplied
as a second unrelated coefficient array. Any two such references for
the same gauge group are equal. This uniqueness concerns perturbative
data, not infrared phases.

\Needspace{8\baselineskip}
\subsection*{G12. Finite normalization measurements}
\label{gloss:G12}
\textbf{Formal:} \code{YangMills.MSbar.Position.Screen}.\newline
\textbf{Source:} \code{MSPositionScreens.lean:54}.

A screen is a finite list of separated profiles and field words,
with a real-linear map from their complex readings to field
endomorphisms. The \code{Prepared} laws require that applying the screen
to the free first variation recovers every allowed infinitesimal
field change. Cutoff compatibility and curvature parity are also
imposed. Thus the screen is a finite determining system for the
normalization frame, rather than a single normalization number.

\Needspace{8\baselineskip}
\subsection*{G13. Signed cubic coupling observable}
\label{gloss:G13}
\textbf{Formal:} \code{YangMills.MSbar.Position.Triangle}.\newline
\textbf{Source:} \code{MSPositionLabels.lean:40}.

The labels are a self-dual cubic curvature field of degree six
and two copies of an anti-self-dual quadratic field of degree four.
They are read at a common cutoff six, with a unit in omitted slots.
The profile is applied after forming the third joint cumulant.
Admissibility requires a real nonnegative, rotation-invariant,
well-separated profile and a negative bare reading $c_T$.
The notation $T$ here is the cubic label, not the stress tensor;
the minus chirality is not charge-conjugation parity.

\Needspace{8\baselineskip}
\subsection*{G14. Independent bare triangle formula}
\label{gloss:G14}
\textbf{Formal:} \code{YangMills.MSbar.Position.bareTriangle}.\newline
\textbf{Source:} \code{MSPositionCubic.lean:85}.

The explicit reference uses the free covariance, orientation and
Lie-bracket color contractions to fix the signed first coefficient.
Both the nonlinear curvature in the insertions and the cubic
interaction contribute. Its integral defines $c_T$ before any
interacting hierarchy is consulted. \code{RawSlopeAgrees} compares
the first connected graph coefficient with this value; the calibration
requires this for every admissible triangle, not just the elected menu.

\Needspace{8\baselineskip}
\subsection*{G15. Formal change of coupling and fields}
\label{gloss:G15}
\textbf{Formal:} \code{YangMills.Matching.Conversion}.\newline
\textbf{Source:} \code{FormalConversion.lean:41}.

The scalar series $f$ satisfies $f_0=0$, $f_1=1$. Field series
$Z$ and $Z^{-1}$ are mutually inverse and preserve homogeneous
weights, the unit, reality, signed axes and cutoff compatibility.
Correlation conversion substitutes $f$ and $Z$ on each input leg
before collecting a coefficient. OPE conversion also applies
$Z^{-1}$ on the output. Omitting that output transformation would
compare coefficients in different field coordinates.

\Needspace{8\baselineskip}
\subsection*{G16. A fixed numerical realization of the field correction}
\label{gloss:G16}
\textbf{Formal:} \code{YangMills.PerturbativeMatching.smoothAdjustment}.\newline
\textbf{Source:} \code{YangMills/PerturbativeMatching/Calibration.lean:50}.

The formal adjustment removes unequal-weight parts of the raw
position frame. A specified smooth cutoff sum realizes all its
Taylor coefficients as a numerical family. Calibration requires
equality with this particular function, fixing even the freedom
invisible to its Taylor series. The group-only choice precedes
the Wilson family. This is different from asking merely for some
smooth realization with the right asymptotic expansion.

\Needspace{8\baselineskip}
\subsection*{G17. Physical position normalization}
\label{gloss:G17}
\textbf{Formal:} \code{YangMills.MSbar.PositionNormalized}.\newline
\textbf{Source:} \code{YangMills/PerturbativeMatching/PositionFrame.lean:36}.

The actual mixed correlators are smeared at physical length $r$
in an unadjusted frame. They approach free readings, satisfy the
finite normalization equations and have an invertible tangent
response for every permitted infinitesimal change. Separate laws
assert local uniqueness, unit insertion, nesting and reality.
These equations use the mixed hierarchy; the physical normalization
is not independently fitted in each character phase.

\Needspace{8\baselineskip}
\subsection*{G18. Measured coupling and field frame}
\label{gloss:G18}
\textbf{Formal:} \code{YangMills.PerturbativeMatching.Measurement}.\newline
\textbf{Source:} \code{YangMills/PerturbativeMatching/Measurement.lean:21}.

One positive coupling and one family of linear field isomorphisms
serve all phases, words and orders on a hierarchy. The coupling
obeys the one-loop leading-log law in the public trace convention.
Field laws include unit insertion, cutoff compatibility, signed-axis
covariance and multilinearity. \code{FrameCalibration} identifies the
coupling exactly with the normalized mixed triangle and the measured
frame with the prescribed adjustment of the position frame.

\Needspace{8\baselineskip}
\subsection*{G19. Meaning of the remainder estimates}
\label{gloss:G19}
\textbf{Formal:} \code{YangMills.Matching.ShapeDerivativeBound}.\newline
\textbf{Source:} \code{YangMills/PerturbativeMatching/Asymptotics.lean:46}.

This predicate requires differentiability through a given order
and a bound on each position derivative of the difference of two
functions. The derivatives are with respect to the separated shape,
not the radius. \code{HasSmearedExpansion} tests all separated
profiles supported in a compact set with bounded derivatives and
gives a little-$o$ remainder at every order. The coefficient norm
is the sum of the absolute polynomial coefficients.

\Needspace{8\baselineskip}
\subsection*{G20. Correlation and stress matching}
\label{gloss:G20}
\textbf{Formal:} \code{YangMills.PerturbativeMatching.CorrelationExpansion}.\newline
\textbf{Source:} \code{YangMills/PerturbativeMatching/Correlations.lean:31}.

The five clauses give smeared expansions, compact
$O(g_X^{p+1})$ bounds, all shape derivatives, short-distance
agreement between any two phase choices, and matching of the
actual physical stress field to the unique Ward series. Constants
have the quantifiers printed in Section~5. In particular the
phase estimate has $r^\gamma$, $0<\gamma<4$, and the stress clause
uses a separate little-$o(g_X^p)$ formulation.

\Needspace{8\baselineskip}
\subsection*{G21. Ordinary OPE characterized by exterior fields}
\label{gloss:G21}
\textbf{Formal:} \code{YangMills.MSbar.Ordinary.Valid}.\newline
\textbf{Source:} \code{MSOrdinaryOPE.lean:69}.

A candidate coefficient array is required to reproduce cluster
expansions of the reference correlations against arbitrary exterior
insertions, with the stated polynomial-logarithmic remainder control.
\code{Determined} says exactly one full array satisfies these laws.
This tests what the fused field does against spectators, not just
its vacuum expectation. It identifies every finite output degree
of the perturbative OPE.

\Needspace{8\baselineskip}
\subsection*{G22. Actual field-valued OPE output}
\label{gloss:G22}
\textbf{Formal:} \code{YangMills.PerturbativeMatching.measuredOPE}.\newline
\textbf{Source:} \code{YangMills/PerturbativeMatching/OperatorProducts.lean:20}.

The physical channel coefficients are multiplied by their
channel fields and included in the space through the chosen
output cutoff. Applying the inverse measured frame gives a free
sector element. \code{HasChannelBasis} requires those included
channels to form a basis. \code{OPEExpansion} then compares the
whole element with the converted graph prediction: smeared
against every output functional, in coefficient norm, and with
shape derivatives uniform over the dual unit ball.

\Needspace{8\baselineskip}
\subsection*{G23. Ward characterization of the stress series}
\label{gloss:G23}
\textbf{Formal:} \code{YangMills.Matching.WardStress.Determined}.\newline
\textbf{Source:} \code{YangMills/PerturbativeMatching/WardIdentities.lean:33}.

The series has the specified leading free stress, symmetry,
signed-axis covariance, vanishing one-point function, conservation
and a translation identity. Translation is expressed by integrating
stress flux against a window isolating the first insertion and
comparing with a derivative of the ordinary correlation test.
Both sides use the same conversion. The signs of the free stress,
the cutoff derivative and integration by parts must be reviewed
together.

\Needspace{8\baselineskip}
\subsection*{G24. Maxwell scaling and all physical radii}
\label{gloss:G24}
\textbf{Formal:} \code{YangMills.Hierarchy.MaxwellScaling}.\newline
\textbf{Source:} \code{YangMills/ScalingLimits.lean:12}.

The scaling object supplies short-distance frames along a
sequence. The matching refinement identifies those sampled frames
with its measured frame and additionally requires
\code{AllRadii.Laws}, \code{AllRadii.FrameBounds} and
\code{FixedFields.Laws}. The full public package therefore has
more content than convergence on a selected sequence of radii.
These ultraviolet requirements do not remove the subsequence
in the large-volume construction.

% END REVIEW_GLOSSARY



\clearpage
\section{Source guide and reproducible provenance}

All Lean paths below are relative to the locally verified November
source tree. The adjacent \code{source-map.json} records its absolute
root, hashes of the cited files, and declaration lines. These are
local review references, not links to an unpublished GitHub repository.
The code is the source for this draft; no informal proof PDF was used.
\code{source-map.json} also records hashes of the verification receipts
and logs used in Section~\ref{sec:receipts}.

\begin{description}[leftmargin=3em,labelwidth=2.6em,align=left]
\item[\src{T}] \textbf{Public target and assembly.}
\code{Comparators/Matching/Challenge.lean};
\code{Comparators/Matching/Solution.lean};
\code{YangMills/PerturbativeMatching.lean}:
\code{MatchedRefinement}, \code{MatchedRun}, \code{MatchedRun.realizes};
\code{YangMills/WilsonProof/Construction.lean}:
\code{refinementOnPublishedRun}, \code{matchedRun};
\code{YangMills/WilsonProof/MeasuredExpansions.lean}.

\item[\src{F}] \textbf{Labels, conversions, and the beta coefficient.}
\code{FormalConversion.lean}:
\code{oneLoopBeta}, \code{Conversion},
\code{Conversion.correlation}, \code{Conversion.ope}.

\item[\src{G}] \textbf{Numerical graph reference.}
\code{YangMills/PerturbativeMatching/GraphCoefficients.lean}:
\code{msbarCorrelationForm}, \code{graphCorrelation};
\code{YangMills/PerturbativeMatching/Reference.lean}:
\code{msbarCoefficients}, \code{Reference}, \code{Reference.scheme_unique};
\code{YangMills/WilsonProof/ReferenceElection.lean}.

\item[\src{D}] \textbf{Graph rules and subtraction.}
\code{YM/MS/Motif/Receipt.lean};
\code{YM/MS/Graphical/Diagram.lean};
\code{YM/MS/Graphical/Enumeration.lean};
\code{YM/MS/Graphical/Lines.lean};
\code{YM/MS/Graphical/Block.lean};
\code{YM/MS/Graphical/PricePin.lean};
\code{YM/MS/Graphical/Panels.lean};
\code{YM/MS/Puff/State.lean}.

\item[\src{R}] \textbf{Regulator-limit theorems.}
\code{YM/Integration/WilsonGraphPublicForestLimits.lean};
\code{YM/Integration/WilsonGraphNamedGraphForestLimits.lean}.

\item[\src{B}] \textbf{Classical input and graph/flow bridges.}
Proof-bearing modules under
\code{YangMills/WilsonProof/NativeSupport/YM/MS/}:
\code{Motif/Pulp.lean}, \code{Motif/Receipt.lean},
\code{Motif/Warp.lean}, \code{Motif/Sync.lean},
\code{Dock/Recipes.lean}, \code{Mach/Land.lean},
\code{Graphical/Lines.lean}, \code{Graphical/Block.lean},
\code{Graphical/Leading.lean}, \code{Graphical/Panels.lean}.
Also \code{YM/MS/Dock/Earn.lean}, \code{YM/MS/Dock/Make.lean},
\code{YM/MS/Dock/Cream.lean}, \code{YM/MS/Swirl/LandLine.lean},
\code{YM/MS/Graphical/CommonGates.lean}, and
\code{YM/Integration/FVA1Physical/Exchange/CalibratedExpansionExternalMSRecursion.lean}.
The prefix matters: several shorter \code{YM/MS/} modules supply
definitions, while the \code{NativeSupport/} modules extend the same
namespaces with the proof theorems cited here.

\item[\src{K}] \textbf{Ghost fields and classical source dictionary.}
\code{YM/MS/Landau.lean}; \code{YM/MS/Properties.lean};
\code{YM/MS/Swirl/LandLine.lean}; \code{YM/MS/Motif/Pulp.lean};
\code{YM/MS/Halo/Task.lean}: \code{Halo.Surf.reach};
\code{YangMills/WilsonProof/NativeSupport/YM/MS/Motif/Beat.lean}:
\code{cape_tick}; the \code{Motif/Pulp.lean} and
\code{Motif/Warp.lean} proof modules under the same prefix.
Our separate file is \code{checks/matching/GhostConventionCheck.lean}
in the verification repository.

\item[\src{S}] \textbf{Poles and finite regulator terms.}
\code{YM/MS/ScaleForest/Bubbles.lean}: \code{TwoPoint.poleK},
\code{TwoPoint.ghost_bill};
\code{YM/MS/ScaleForest/SternGroup.lean}: \code{Triangle.ghost_input};
\code{YM/MS/ScaleForest/TestScale.lean}: \code{Triangle.premature},
\code{Triangle.DScaleJet};
\code{YM/MS/ScaleForest/PageOne.lean}: \code{simple_finite},
\code{simple_pole_rule}; \code{YM/MS/ScaleForest/Entry.lean}:
\code{final_including}; \code{YM/MS/ScaleForest/QuarticRule.lean}:
\code{sums_angles}; \code{YM/MS/ScaleForest/TwoVertices.lean}:
\code{twoVertex_laurent}; \code{YM/MS/ScaleForest/HeatScale.lean}:
\code{halfP_components}; \code{YM/MS/Image/Seat.lean}:
\code{Outlook.four}; \code{Pearl/Broach.lean} under the
\code{NativeSupport/YM/MS/} prefix in [B]: \code{Pearl.chosen}.
Our files under \code{checks/matching/} are \code{CountertermAudit.lean}
(type and axiom inspection) and \code{GhostLoopBenchmark.lean}
(reference calculation and graph comparisons).

\item[\src{N}] \textbf{Normalization and measured fields.}
\code{MSPositionScreens.lean}; \code{MSPositionPrescription.lean};
\code{MSConversion.lean};
\code{YangMills/PerturbativeMatching/Calibration.lean};
\code{YangMills/PerturbativeMatching/PositionFrame.lean};
\code{YangMills/PerturbativeMatching/Measurement.lean}.

\item[\src{C}] \textbf{The cubic diagnostic.}
\code{MSPositionLabels.lean}:
\code{cubicValue}, \code{quadraticValue}, \code{Triangle},
\code{Triangle.RawSlopeAgrees};
\code{MSPositionCubic.lean}: \code{bareTriangle}.

\item[\src{A}] \textbf{Exact asymptotic requirements.}
\code{YangMills/PerturbativeMatching/Asymptotics.lean};
\code{YangMills/PerturbativeMatching/Correlations.lean};
\code{YangMills/PerturbativeMatching/OperatorProducts.lean}.

\item[\src{V}] \textbf{Lattice and volume connections.}
\code{YangMills/PerturbativeMatching/NativeLattice.lean};
\code{YangMills/PerturbativeMatching/VolumeLimits.lean}:
\code{SmoothMeshLimit}, \code{HasTestLimits}, \code{RepresentsTestLimits}.

\item[\src{P}] \textbf{Native expansion and its public identification.}
\code{YM/SignedWilsonExpansion/Result/Orders.lean}:
\code{LinePrices}, \code{ConferencePowers}, \code{pricedLines};
\code{YangMills/WilsonProof/NativeSupport/ContinuumConstruction/StandaloneCorrelation.lean};
\code{YM/Integration/WilsonGraphMeasuredCompactCorrelation.lean}:
\code{publicMixedPolynomial}, \code{publicMixedCompact}.

\item[\src{E}] \textbf{Equality of the calibrated graph and Wilson coefficients.}
\code{YM/MS/Graphical/CommonGates.lean};
\code{YM/Integration/FVA1Physical/Exchange/CalibratedExpansionCommonCalibratedGates.lean}:
\code{original_calibrated_gate_of_team},
\code{original_filed_calibrated_gates}.

\item[\src{J}] \textbf{Derivatives and the correlation assembly.}
\code{YM/Integration/WilsonGraphCompactCoordinateJets.lean};
\code{YM/Integration/WilsonGraphAppointedStationaryPrices.lean};
\code{YM/Integration/WilsonGraphAppointedMixedShapeJets.lean};
\code{YM/Integration/WilsonGraphAppointedPhysicalShapePrice.lean};
\code{YM/Integration/WilsonGraphPublishedShapeEstimate.lean};
\code{YM/Integration/WilsonGraphNamedCorrelationEstimates.lean};
\code{YangMills/WilsonProof/NativeSupport/CorrelationMatching/Construction/Estimates.lean}.

\item[\src{H}] \textbf{Phase proximity.}
\code{YM/OPE/Short/Gap.lean};
\code{YM/PhysMS/Foundation/Readout.lean};
\code{YM/PhysMS/Foundation/Estimates.lean};
\code{YM/PhysMS/Frame/Phases.lean}.

\item[\src{O}] \textbf{Reference OPE and uniqueness.}
\code{MSOrdinaryOPE.lean};
\code{YangMills/WilsonProof/NativeSupport/IssuedMaxwell/GraphOnlyPrice.lean};
\code{YangMills/WilsonProof/NativeSupport/IssuedMaxwell/GraphOrdinary.lean};
\code{YangMills/WilsonProof/NativeSupport/ContinuumConstruction/GroupOrdinaryTarget.lean}.

\item[\src{Q}] \textbf{Full physical OPE output estimates.}
\code{YM/Integration/WilsonFlashOnThePublishedBasisRows.lean};
\code{YM/Integration/WilsonFlashRoundedOnTheElectedWindow.lean};
\code{YM/Integration/WilsonFlashOnBothPublishedRows.lean};
\code{YangMills/WilsonProof/NativeSupport/CommissionedPublic/AppointedRoundedReader.lean};
\code{YM/Integration/WilsonGraphAppointedOrdinaryTestedExpansion.lean};
\code{YangMills/WilsonProof/NativeSupport/WilsonFlashScoredOrdinaryEstimates.lean}.

\item[\src{W}] \textbf{Stress normalization and matching.}
\code{YangMills/PerturbativeMatching/WardIdentities.lean}:
\code{WardStress.Determined};
\code{YangMills/PerturbativeMatching/Correlations.lean}:
\code{CorrelationExpansion.stressExpansion};
\code{YangMills/WilsonProof/NativeSupport/WilsonFlashScoredGroupReference.lean}.

\item[\src{Z}] \textbf{Verification records.}
Under \code{work/ym-nov4-verification/evidence}:
\code{matching-comparator.json}, \code{matching-comparator.log},
\code{patched-matching-comparator.json}, \code{patched-matching-comparator.log},
\code{matching-axioms.json}, \code{matching-axioms.log}.
The word ``patched'' in the record names identifies the repeat with the
newer upstream checker; it does not indicate an edited mathematical
Matching target.
\end{description}

\clearpage
\section{Explicit finite ghost-loop calculation}
\label{app:finite-ghost}

This appendix computes the finite constants from the submission's
Gaussian and Mellin prescription. It provides a more direct comparison
than first identifying the entire graph with a dressed two-propagator
integral. The resulting answer agrees with the independent
$\MS$ expression in~\eqref{eq:ghost-loop-benchmark}, including both
finite constants. \emph{The full comparison is now proved in Lean and
has passed Comparator.} The file
\code{checks/matching/GhostLoopFiniteEvaluation.lean}
proves the ordinary source evaluation, the scalar radial integrals,
the polynomial identities, and the exact parameter integrals used
below. The subsequent modules connect the actual sector contractions,
\code{closeG} integral and regulator derivative to this calculation.
They then determine the finite amplitude from its scale derivative
and ultraviolet boundary and prove its infrared limit. We first give
the analytic calculation, then explain the precise route taken by
the completed Lean proof.

\subsection{The force before the scale integration}

Let $C=C_A\delta^{ab}$, and initially choose units with $p^2=1$.
Write $u=\delta_{\mu\nu}$ and $v=p_\mu p_\nu$. All formulas retain
the sign of~\eqref{eq:ghost-loop-integral} and omit the overall
coupling factor $g^2$. Put $D=4+2z$, so that $z=-\varepsilon$.
The force $W(z)$ is \code{TwoPoint.ratePair G true true}, evaluated
on the external two-gauge-field word at momentum $\lambda p$.

At $z=0$, \code{submitted_force_four} identifies the actual submitted
force, for every diagonal tensor component, with $-8C$ times the
integral of \code{SquareShift.heatValue} over $0\le t\le1$.
This uses \code{log_rate}, \code{ghostTensor}, \code{from_integral},
and \code{payment_line}. The factor includes both ghost orientations
and the source's contraction conventions; it is not an additional
choice of a loop sign. The polynomial before Gaussian integration is
\[
 q^2(q+\lambda p)^2(q+\lambda p)_\mu q_\nu.
\]
The two powers of squared momentum come from the differentiated heat
lines. They disappear only after the scale integrations; omitting them
at this stage would change the calculation.

Complete the square and change variables by
\[
 x=\frac{t^2}{1+t^2},\qquad r=1-x,\qquad
 0\le x\le\tfrac12,\qquad \nu=1+t^2=r^{-1}.
\]
The Gaussian covariance is $\sigma=r/2$, and the two momentum shifts
are $-x\lambda p$ and $r\lambda p$. The continued expression obtained
from the split-radial integral is
\begin{equation}\label{eq:continued-force-explicit}
 W(z;\lambda p)=-\frac{4C}{(4\pi)^{2+z}}
  \int_0^{1/2}x r^{z-1}e^{-x\lambda^2}
   \bigl(a_D+b_D\lambda^2+c_D\lambda^4+d_D\lambda^6\bigr)\,dx.
\end{equation}
The derivation of its dimensional dependence is given next; agreement
with the ordinary four-dimensional expression alone would not establish
this equation or its regulator derivative.

\subsection{Why the split radial integral gives these moments}

In the physical external sector, \code{closeG} integrates a
four-dimensional momentum $q$ and an auxiliary radial coordinate $R$.
Its factors \code{gage 0}, \code{pruning 1}, and \code{feeD} give,
for an even smooth radial function, the normalized operation
\begin{equation}\label{eq:radial-functional}
 \mathcal D_z f=-\frac1{\Gamma(1+z)}
          \int_0^\infty R^{2z} f'(R)\,dR,
\end{equation}
together with $(2\pi)^{-4}(4\pi)^{-z}$ and the ordinary $q$ integral.
Indeed, \code{gage 0} is $-1$ and
\code{pruning 1} applied to $f$ is $f'(R)/R$ for $R>0$; the latter follows by
integrating the second derivative and using $f'(0)=0$.
The defining exponent in \code{closeG} is $R^{2z+1}$.

There is no separate continued polarization trace for this graph.
Our \code{ghost_vector_line_zero} proves that the selected ghost line
vanishes on two vector legs in every test sector.
\code{ghost_yeast_zero} and \code{ghost_innerTrace_no_explicit_z}
therefore remove the explicit $2(z-m)$ times \code{yeast} term in
\code{innerTrace}. Our \code{ghost_loopCarpet_regulator} and
\code{ghost_closeG_physical} also prove that the actual closure uses
the fixed radial function at $z=0$, weighted by $R^{2z+1}$ and
\code{feeD}. Regulator dependence still enters through
\eqref{eq:radial-functional}; removing the polarization term does not
make $W$ independent of $z$.

For $\operatorname{Re}z>-1$ the three radial moments are
\begin{align}
 \mathcal D_z(e^{-\nu R^2})&=\nu^{-z},\nonumber\\
 \mathcal D_z(R^2e^{-\nu R^2})&=z\nu^{-z-1},\label{eq:three-radial-moments}\\
 \mathcal D_z(R^4e^{-\nu R^2})&=z(z+1)\nu^{-z-2}.\nonumber
\end{align}
They follow by differentiating the displayed functions, substituting
$w=\nu R^2$, and applying the Gamma recurrence. Our three
\code{normalized_radial} lemmas prove these scalar integral
expressions for real $z>-1$, including a neighborhood of zero.
Their identification with the actual source integral is proved in
\code{GhostLoopContinuation.lean}; the source analytic germ then fixes
the complex jet in \code{GhostLoopJet.lean}.

Here is an explicit check that this continuation gives the expected
Gaussian polynomial. Let $M_D$ denote the moment of
\[
 |X+ap|^2|X+bp|^2(X+bp)_\mu(X+ap)_\nu
\]
for integer $D$ and an isotropic Gaussian of covariance $\sigma$,
with $p^2=Q$, then continue the resulting polynomial in $D$. The polynomial \code{momentD} is obtained by Wick
contraction. Its $D=4$ value is exactly the submission's
\code{pairValue}, as proved by \code{moment_four}.
Before integrating $R$, the ordinary Gaussian moment is
$M_4+R^2E_1+R^4E_2$, where
\begin{align*}
 E_1={}&u\bigl(12\sigma^2+(a^2+b^2)Q\sigma\bigr)\\
 &+v\bigl(2(a+b)^2\sigma+ab(8\sigma+(a^2+b^2)Q)\bigr),\\
 E_2={}&\sigma u+abv.
\end{align*}
The factors of $R^2$ enter only the two squared momenta, because
the external indices and momenta lie in the physical plane.
Equation~\eqref{eq:three-radial-moments}, with $\nu^{-1}=2\sigma$,
then gives
\[
 M_{4+2z}=M_4+2z\sigma E_1+4z(z+1)\sigma^2E_2.
\]
\code{split_radial_polynomial} proves this polynomial identity.
Thus the dimension-dependent moment is obtained by evaluating the
source's radial prescription, rather than interpolating its values
at integer dimensions.

For this Gaussian integrand the required interchanges have elementary
bounds. Near $R=0$ the differentiated even polynomial times a Gaussian
is $O(R)$, so $R^{2z}f'(R)$ and its first $z$ derivative are integrable
uniformly on compact subsets of $\operatorname{Re}z>-1$.
At infinity the Gaussian controls every polynomial and logarithm.
The $t$ integration is over $[0,1]$ with $\nu\ge1$ and a $t^3$
factor. These bounds justify the Gaussian and radial integrations
and differentiation at zero in the analytic calculation.
The corresponding source contraction, measure maps and integral
exchanges are formalized in \code{GhostLoopTrace.lean},
\code{GhostLoopRadial.lean} and \code{GhostLoopContinuation.lean}.
The first regulator derivative is obtained from the resulting compact
parameter integral in \code{GhostLoopJet.lean}.

\subsection{The coefficients and the local regulator jet}

Expand $M_D$ using $Q=1$, $a=-x\lambda$, $b=r\lambda$,
and $\sigma=r/2$. The coefficients in
\eqref{eq:continued-force-explicit} at $D=4$, and their derivatives
with respect to $z$ through $D=4+2z$, are as follows. Each column
is the coefficient of the indicated tensor structure.
\begin{center}
\begin{tabular}{c|c|c}
 & $u$ & $v$\\ \hline
$a$ & $6r^3$ & $0$\\
$b$ & $\frac{r^2}{2}(3-8x+8x^2)$ & $r^2(3-20x+20x^2)$\\
$c$ & $\frac{x^2r^3}{2}$ & $-xr^2(3-10x+10x^2)$\\
$d$ & $0$ & $-x^3r^3$\\
$a'$ & $\frac72r^3$ & $0$\\
$b'$ & $\frac{r^2}{2}(1-2x+2x^2)$ & $r^2(1-9x+9x^2)$
\end{tabular}
\end{center}
The declarations \code{coeffA}, \code{coeffB}, \code{coeffC}, and
\code{coeffD} give the polynomials at general $D$.
\code{moment_expansion} verifies the expansion and
\code{local_taylor_derivative} verifies $b'-xa'$.

The quadratic Taylor coefficient of the exponential times the
polynomial is $b_D-xa_D$. Since
$T(z)=-\tfrac12\operatorname{Tad}_4W(z)$, differentiating
\eqref{eq:continued-force-explicit} gives
\begin{equation}\label{eq:ghost-jet-explicit}
 T'(0;p)=C R_{\mu\nu}(p)
                 \left(\frac{11}{24}+\log(2\pi)\right).
\end{equation}
For example the $u$ channel at $p^2=1$, after factoring out
$C/(16\pi^2)$, is
\[
 \int_0^{1/2}\!\left\{J(x)+W_\ell(x)
                 [\log(1-x)-\log(4\pi)]\right\}dx,
\]
where $J=xr(1-9xr)$ and $W_\ell=xr(3-20xr)$.
The three needed integrals are
\[
 \int_0^{1/2}J=-\frac1{15},\qquad
 \int_0^{1/2}W_\ell=-\frac1{12},\qquad
 \int_0^{1/2}W_\ell\log(1-x)
                  =\frac{151}{1440}-\frac{\log2}{12}.
\]
Our \code{regulator_jet_parameter} proves their combination is
$11/288+\log(2\pi)/12$. The $v$ channel is twice the $u$ channel.
The Lean theorem \code{GhostJet.submitted_regulator_jet_test}
proves \code{SubmittedRegulatorJetTest}, namely
\eqref{eq:ghost-jet-explicit} for the actual
\code{ruby (highTensor 4 ratePair) 1}. This is a source-level equality,
proved using the continued radial integral and the analytic germ;
it is not a hypothesis of the final finite-part theorem.

\subsection{The Mellin finite part and removal of the infrared cutoff}

For this external word, the scale integral has weight
$\lambda^{-3-2z}d\lambda$. At a fixed $x$, its continued value is
\begin{equation}\label{eq:ghost-mellin-polynomial}
 \frac12\left[
 a_Dx^{1+z}\Gamma(-1-z)+b_Dx^z\Gamma(-z)
 +c_Dx^{-1+z}\Gamma(1-z)+d_Dx^{-2+z}\Gamma(2-z)\right].
\end{equation}
This is \emph{not} an ordinary convergent integral at $z=0$.
One subtracts the constant and quadratic Taylor terms at
$\lambda=0$, integrates the remainder, and adds the corresponding
meromorphic elementary integrals. This is precisely the decomposition
into \code{ordinary} and $z^{-1}$ times \code{highTensor} used in
\code{PageOne.one_loop_germ}.
With the four coefficients held at their $z=0$ values, the finite
part of~\eqref{eq:ghost-mellin-polynomial} is
\[
 \tfrac12\left[
 ax(\gamma_E-1+\log x)-b(\gamma_E+\log x)+c/x+d/x^2\right].
\]
The derivatives $a',b'$ also contribute through the two pole terms;
they must not be discarded.

The infrared removal can be done in this same representation.
For fixed $p\ne0$, equation~\eqref{eq:continued-force-explicit}
and the table give $W(0;\lambda p)=O(\lambda^{-2})$ as
$\lambda\to\infty$. For instance the term with $b\lambda^2$
is bounded by a constant times
$\lambda^2\int_0^\infty x e^{-c\lambda^2x}dx=O(\lambda^{-2})$;
the other terms obey the same or a stronger bound. Here $c=p^2>0$.
The bound also holds uniformly for the regulator in a small compact
neighborhood of zero, since $1/2\le r\le1$.
Consequently the upper Mellin tail is $O(\rho^{-4})$ at $z=0$.

More explicitly, write $A_0=W(0;0)$ and let $A_2(p)$ be the
coefficient of $\lambda^2$ in $W(0;\lambda p)$. Then
$A_2=-2T(0)$, and the source's regular part is
\[
 \mathcal R_{\log\rho}(p)=
  \int_0^\rho\!\frac{W(0;\lambda p)-A_0-\lambda^2A_2(p)}{\lambda^3}
          \,d\lambda-\frac{A_0}{2\rho^2}.
\]
The logarithm from the subtracted quadratic term is canceled by
$-2\log\rho\,T(0)$ in~\eqref{eq:retained-ghost-finite}.
After cancellation, differentiating the finite expression with
respect to $\log\rho$ gives $\rho^{-2}W(0;\rho p)$.
Its integrable tail proves convergence and the stated error estimate.
This also explains why pointwise convergence of a propagator was
insufficient on its own: the Taylor subtractions and their finite
regulator term must be carried through the limit.

At $\bar\mu=1$ and $p^2=1$, combining the Mellin finite part,
$T'(0)$, and $y=(\gamma_E-\log(4\pi))/2$ cancels $\gamma_E$
and $\log(4\pi)$. The remaining integrand, before $C/(16\pi^2)$,
is
\[
 \frac{2x}{r}\left[xa+b'-xa'-c/x-d/x^2
                       +(b-xa)\{\log x+\log r\}\right].
\]
\code{finite_integrand_reduction} identifies it with
\[
 u\{P_\delta+W_\ell\log(xr)\}
   +v\{P_p+2W_\ell\log(xr)\},
\]
where $P_\delta=xr(1+2xr)$ and $P_p=2xr(4-18xr)$.
Each combined integrand is symmetric under $x\mapsto1-x$.
Exact elementary integrals give
\[
 \int_0^1P_\delta=\frac7{30},\quad
 \int_0^1P_p=\frac2{15},\quad
 \int_0^1W_\ell\log x=\frac{19}{180}.
\]
Thus the two half-interval values are
\begin{equation}\label{eq:finite-ghost-constants}
 \frac7{60}+\frac{19}{180}=\frac29
       =\frac1{12}\frac83,\qquad
 \frac1{15}+\frac{19}{90}=\frac5{18}
       =\frac16\frac53.
\end{equation}
These are exact Lean integral equalities
\code{finite_delta_half} and \code{finite_pair_half}, not numerical
quadrature. The endpoint logarithms are handled through continuous
$x^n\log x$ and an endpoint-safe primitive.

Momentum rescaling restores $p^2$ and changes the logarithm to
$-\log p^2$. Adding the already proved renormalization-scale change
then gives the full analytic answer
\begin{equation}\label{eq:finite-ghost-answer}
 B^{ab}_{\mu\nu,\MS}(p)=\frac{C_A\delta^{ab}}{16\pi^2}
 \left[
 \frac{\delta_{\mu\nu}p^2}{12}
       \left(\log\frac{\bar\mu^2}{p^2}+\frac83\right)
 +\frac{p_\mu p_\nu}{6}
       \left(\log\frac{\bar\mu^2}{p^2}+\frac53\right)
 \right],\qquad p^2>0.
\end{equation}
It agrees with the independent Gamma-function benchmark and the
finite constants in the cited calculation of Rodigast.

\subsection{How the Lean proof fixes the finite amplitude}
\label{sec:finite-proof}

The formal proof does not require exchanging a regulator derivative
with the full improper scale integral. It first evaluates the continued
force and its local Taylor coefficient, then solves the scale equation
at $z=0$ with the source's ultraviolet boundary condition. This route
keeps the subtraction constant visible and proves an equality at every
retained infrared cutoff, before taking its limit.

For a diagonal component at $p^2=1$, put $v=p_i^2$ and $X=\rho^2$.
The actual source force is
\[
 W(0;\rho p)=-\frac{4C}{16\pi^2} f_v(X),
\]
where the compact parameter integral evaluates to
\begin{align*}
 f_v(X)={}&e^{-X/2}\left[
 \frac{2-8v}{X^2}+\frac{-1/2-v}{X}-\frac5{32}+\frac v2
 +\left(-\frac1{64}+\frac v8\right)X+\frac v{64}X^2\right]\\
 &+\frac{-2+8v}{X^2}+\frac{3/2-3v}{X}.
\end{align*}
The theorem \code{GhostInfrared.submitted_force_unit} proves this
identity for the submitted \code{ratePair}; it is not a new definition
of that force.

Define $k_v=(1+2v)/48$ and the explicit primitive
\begin{align*}
 P_v(X)={}&e^{-X/2}\left[
 \frac{-1+4v}{3X^3}+\frac{5/24-v/12}{X^2}
 +\frac{-5/192-5v/24}{X}-\frac v{64}\right]\\
 &-\frac{-1+4v}{3X^3}+\frac{-3/8+3v/4}{X^2}
       +k_v E_1(X/2),
\end{align*}
where $E_1$ is the usual exponential tail. In Lean it is represented by
\[
 E_1(X)=-\gamma_E-\int_0^X\log t\,e^{-t}\,dt-e^{-X}\log X
 \qquad(X>0).
\]
The integral over $(0,\infty)$ of $\log t\,e^{-t}$ is proved to be
$-\gamma_E$ using the derivative of the Gamma function. The derivative
and both endpoint limits of $P_v$ are then established:
\begin{align*}
 P_v'(X)&=\frac{f_v(X)}{2X^2},\qquad
 \lim_{X\to\infty}P_v(X)=0,\\
 \lim_{X\downarrow0}\left(P_v(X)+\frac{11}{64X}+k_v\log X\right)
 &=c_v:=\frac{53}{1152}+\frac{29v}{576}
                     +k_v(\log2-\gamma_E).
\end{align*}
The small-$X$ limit uses a controlled Taylor remainder for the
exponential; none of these constants is inferred from numerical data.

On the source side, \code{GhostFlow.finite_derivative} differentiates
\eqref{eq:retained-ghost-finite} with respect to $s=\log\rho$.
The quadratic Taylor subtraction cancels the explicit logarithmic
term, leaving $e^{-2s}W(0;e^sp)$. The separate theorem
\code{ordinary_boundary} proves the ultraviolet normalization
\[
 \mathcal R_s(p)-\frac{11C}{16(16\pi^2)}e^{-2s}\longrightarrow0
 \qquad(s\to-\infty).
\]
Together with the small-$X$ expansion, this fixes the integration
constant. Equality of derivatives and this boundary value give the
exact identity \code{GhostUnit.finite_unit_exact}:
\begin{equation}\label{eq:finite-cutoff-primitive}
 B_s(y;p)=-\frac{4C}{16\pi^2}P_v(e^{2s})
          +\frac{4C}{16\pi^2}c_v+T'(0;p)-2yT(0;p).
\end{equation}
The large-$X$ limit removes the first term. Inserting the proved
regulator jet and the prescribed $\MS$ scale gives the two constants
in~\eqref{eq:finite-ghost-answer}.

The remaining quantifiers are also proved. Rotational covariance
recovers all tensor components from the two diagonal probes. The exact
scaling law
\[
 B_s(y;xp)=x^2 B_{s+\log x}(y-\log x;p),\qquad x>0,
\]
reduces any $p^2>0$ to unit momentum and produces $\log(\bar\mu^2/p^2)$.
The color contraction is identified with $C_A\delta^{ab}$ using the
submission's simple-gauge-group identity. Finally the previously
proved finite-reader bridge replaces $B_s$ with the actual
\code{Outlook.four} value in the frozen target.

\subsection{Proof files and verification records}

All proof files below are in \code{checks/matching/}. The names in the
second column are declarations to start reading, with namespace
prefixes omitted where the file supplies them.
\begin{longtable}{@{}p{0.31\linewidth}p{0.65\linewidth}@{}}
\toprule
File & Role in the completed proof\\\midrule\endhead
\code{GhostLoopBenchmark.lean} & Independent finite reference, actual
reader, pole and scale checks; definition of
\code{SubmittedFinitePartTest}.\\
\code{GhostLoopTrace.lean} & Actual ghost contractions in all auxiliary
sectors: \code{trace_value} and \code{inside_loop}.\\
\code{GhostLoopRadial.lean} & Sector lifts, Gaussian moments and justified
integral exchanges: \code{actual_carpet_evaluated}.\\
\code{GhostLoopContinuation.lean} & Continued radial measure and actual
source force: \code{actual_force_real}, \code{actual_force_scaled}.\\
\code{GhostLoopJet.lean} & First regulator derivative of the actual local
term: \code{submitted_regulator_jet_test}.\\
\code{GhostLoopInfrared.lean} & Explicit physical force:
\code{submitted_force_unit}.\\
\code{GhostLoopMellin.lean} & Logarithmic exponential integral, tail
derivative and endpoint limits.\\
\code{GhostLoopPrimitive.lean} & Primitive derivative and the two
endpoint limits \code{primitive_small}, \code{primitive_large}.\\
\code{GhostLoopFlow.lean} & Submitted scale derivative and ultraviolet
normalization: \code{finite_derivative}, \code{ordinary_boundary}.\\
\code{GhostLoopUnit.lean} & Equality at every cutoff and its limit:
\code{finite_unit_exact}, \code{finite_unit_limit}.\\
\code{GhostLoopTensor.lean} & Tensor covariance and momentum scaling:
\code{unit_tensor}, \code{finite_scale_word}.\\
\code{GhostLoopSolution.lean} & All nonzero momenta and the complete
acceptance theorem \code{GhostSolution.acceptance}.\\
\bottomrule
\end{longtable}
The earlier \code{GhostConventionCheck.lean} and
\code{GhostLoopFiniteEvaluation.lean} supply the vertex conventions
and elementary integrals; \code{GhostFiniteAudit.lean} prints the
final theorem and seven principal axiom reports.

On 13 November, all 15 review modules compiled in a fresh overlay
against the pinned, already-built November source. The seven principal
reports listed only \code{propext}, \code{Classical.choice} and
\code{Quot.sound}. One review-code warning suggested a simpler tactic
sequence; there were no build errors. The source files and logs are
hashed in \code{evidence/matching/ghost-finite-build.json}.

\Needspace{6\baselineskip}
The subsequent Comparator run checked
\begin{quote}
\code{YMReview.NovemberMatching.GhostSolution.acceptance (G : YM.Gauge)}\par
\code{: YMReview.NovemberMatching.GhostLoop.SubmittedFinitePartTest G}.
\end{quote}
The reference imports the benchmark, not the solution. There are no
definition holes. Comparator exported and rechecked the proof,
reported that Lean's default kernel accepted it, and finished with
\texttt{Your solution is okay!} in about 147 seconds. All 37,673
submission files were checked against the pinned manifest before and
after the run, as were the frozen review inputs. The run uses Comparator
revision \code{d03acab154d269c06e60e4de7e4cc85deebff94b} and the Lean
4.34.0 checking kernel, exporting the proof from Lean 4.32.0.
The full log, configuration and stability records are retained under
\code{evidence/matching/ghost-finite-comparator/}; reproduction commands
are in \code{docs/matching/ghost-finite-check.md}.

This was a macOS development-mode run with the upstream
\code{fake-landrun.sh} launcher. It did not provide Linux sandbox
isolation, the optional second kernel, or a Palomar check. The reference's
intentional theorem placeholder produces a \code{sorry} warning while
building the challenge; the exported solution and the audited final
theorem do not depend on \code{sorryAx}. The log also replays warnings
from the already-built submitted modules.

\subsection{Scope of the conclusion}

The full finite-value test passes without a finite-part hypothesis.
It checks a concrete contribution beyond its pole or the coefficient
of its scale logarithm, including the finite effect of the regulator
derivative. It does not validate every graph, composite insertion,
or complete gauge-invariant correlation. The conventional integral's
reduction to the independent Gamma expression in Section~3.5 is still
a human calculation; the formal theorem identifies the actual source
graph with the explicit finite expression stated there. Thus it is
substantial evidence for the finite subtraction conventions in this
example, while the physical coefficient tests I2--I4 and the wider
matching review remain separate work.

\end{document}
